\documentclass[a4paper,11pt]{memoir}

\usepackage[utf8]{inputenc}
\usepackage[T1]{fontenc}
\usepackage[svgnames]{xcolor}
\usepackage{libertine}
\usepackage[frenchb]{babel}
% TODO \usepackage[french]{babelbib}
\usepackage{mathpartir}
\usepackage{syntax}
\usepackage{amssymb}
\usepackage{tikz}
\usepackage{amsmath}
\usepackage{float}
\usepackage{url}
\usepackage{xspace}
\usepackage{graphicx}
\usepackage[shortlabels]{enumitem}
\usepackage[backref=page]{hyperref}
\usepackage{links}
\usepackage[a4paper]{geometry}
\usepackage{amsthm}
\usepackage{stmaryrd}
\usepackage[normalem]{ulem}
\usepackage{wasysym}
\usepackage{multirow}
\usepackage{thmtools}
\usepackage{gramster}
\usepackage{switch}
\usepackage{pelagiccod}
\usepackage{conteq}

\newcommand{\irule}[3]{ \inferrule*[right=(#1)]{#2}{#3} }
\def\iaxiom#1#2{\irule{#1}{ }{#2}}

% http://tex.stackexchange.com/questions/3964/mathematical-symbol-for-disjoint-set-union
\makeatletter
\def\moverlay{\mathpalette\mov@rlay}
\def\mov@rlay#1#2{\leavevmode\vtop{%
   \baselineskip\z@skip \lineskiplimit-\maxdimen
   \ialign{\hfil$\m@th#1##$\hfil\cr#2\crcr}}}
\newcommand{\charfusion}[3][\mathord]{
    #1{\ifx#1\mathop\vphantom{#2}\fi
        \mathpalette\mov@rlay{#2\cr#3}
      }
    \ifx#1\mathop\expandafter\displaylimits\fi}
\makeatother

\newcommand{\dcup}{\charfusion[\mathbin]{\cup}{+}}
\newcommand{\bigdcup}{\charfusion[\mathbin]{\bigcup}{+}}

\def\opbin{\boxplus\xspace}
\def\opun{\boxminus\xspace}

% Types
\newcommand{\tInt}{\textsc{Int}\xspace}
\newcommand{\tFloat}{\textsc{Float}\xspace}
\newcommand{\tUnit}{\textsc{Unit}\xspace}

\def\structRec#1#2{\textsc{rec}(\&#1) #2} % TODO

\newcommand{\tFV}[1]{\textrm{FreeVars}(#1)}

\def\opxor{\mathbin{\char`\^}}

% Constant
\newcommand{\cNil}{\textsc{Nil}}

% Program & language names
\newcommand{\ctonewspeak}{\textsc{c2newspeak}}
\newcommand{\cpp}{\textsc{cpp}}
\newcommand{\newspeak}{\textsc{Newspeak}}

% Newspeak
\newcommand{\npkIf}[3]{\textsc{If}(#1) \{ #2 \} \textsc{Else} \{ #3 \}}
\newcommand{\npkDoWith}[2]{\textsc{Do}\{ #1 \}\textsc{With}~#2:}
\newcommand{\npkWhile}[1]{\textsc{While}(1)\{#1\}}
\newcommand{\npkGoto}[1]{\textsc{Goto}~#1}
\newcommand{\npkReturn}[1]{\textsc{Return}~#1}

% cML
\def\iPass{\textsc{Pass}\xspace}
\def\iIf#1#2#3{\npkIf{#1}{#2}{#3}}
\def\iWhile#1#2{\textsc{While}(#1)\{#2\}}
\def\iReturn#1{\textsc{Return}(#1)}
\def\iDecl#1#2#3{\textsc{Decl}~#1=#2~\textsc{in}\{#3\}}

\def\iThen#1#2{\textsc{If}(#1) \{ #2 \}}

\def\eNull{\textsc{Null}\xspace}
\def\eUnit{(\,)\xspace}

\def\vRet{\underline{R}}

% sets pour semantique
\def\sId{\textsc{Id}\xspace}
\def\sMem{\textsc{Mem}\xspace}
\def\sAddr{\textsc{Addr}\xspace}
\def\sInt{\textsc{Int}\xspace}
\def\sFloat{\textsc{Float}\xspace}
\def\sPhi{\Phi}

% Memory state (instructions)
% #1 : memory
% #2 : instruction
\def\msi#1#2{ \langle #2, #1 \rangle }

% Reduction (expressions)
\def\mm#1#2#3#4{ \msi{#1}{#2} → \msi{#3}{#4} }
\def\mms#1#2{\mm{m}{#1}{m}{#2}}

% Reduction (instructions)
\def\mmi#1#2#3#4{ \mm{#1}{#2}{#3}{#4} }

% Reduction (left-values)
\def\mmlv#1#2#3#4{ \mm{#1}{#2}{#3}{#4} }

% Typage sém
\def\semtyp#1#2#3{#1 ⊧ #2 : #3}

% Compatibilité Γ / m
\def\mcomp#1#2{#1 ⊧ #2}

% Reduction contexts
\def\ctxOp#1#2{#1~\opbin~#2}
\def\ctxEmpty{\bullet\xspace}
\def\ctxUnOp#1{\opun~#1}
\def\ctxLvDeref#1{*#1}
\def\ctxLvField#1#2{#1.#2}
\def\ctxLvIndex#1#2{#1[#2]}
\def\ctxSet#1#2{#1 ← #2}

% Context substitution
\def\ctxSub#1#2{#1\llparenthesis#2\rrparenthesis}
\def\ctxSubL#1#2{\ctxSub{#1}{#2}_L}
\def\ctxSubI#1#2{\ctxSub{#1}{#2}_I}

% Qualifs
\def\qKernel{\textsc{Kernel}\xspace}
\def\qUser{\textsc{User}\xspace}
\def\vTainted#1{\widehat{\diamondsuit}~#1}
\def\eTaint#1{\diamondsuit~#1}

\def\ptrK#1{#1~*}
\def\ptrU#1{#1~@}

% Eval de phrase
\def\ph#1#2#3{#1 ⊩ #2 → #3}

\def\phstar#1#2#3{#1 ⊩ #2 →^* #3}

\def\phprog#1#2{⊩ #1 →^* #2}

% typage
\def\ty#1#2#3{#1 ⊢ #2 : #3}
\def\tyi#1#2{#1 ⊢ #2}
\def\typh#1#2#3{#1 ⊢ #2 → #3}
\def\typhstar#1#2#3{#1 ⊢ #2 →^* #3}

\def\sVal{\textsc{Val}}
\def\sByte{\textsc{Byte}}

\def\powerset#1{\mathcal P \left({#1}\right)}

% Functions with placeholders
\def\phx#1{\ensuremath{#1{\cdot}}}
\def\phxx#1{\ensuremath{#1{\cdot}{\cdot}}}

% Progs bien formés
\def\WF#1{\textsc{wf}(#1)}
\def\NORET#1{\textsc{NoRet}(#1)}

\newcommand{\insertcode}[1]{\input{gen/#1.pyg.tex}}

% \chapter* + add to toc
%\newcommand{\chapterstar}[1]{\chapter*{#1} \addcontentsline{toc}{chapter}{#1}}

% Register names
\def\defreg#1{\expandafter\def\csname#1\endcsname{\textsc{#1}\xspace}}

\defreg{eax}
\defreg{ebx}
\defreg{ecx}
\defreg{edx}
\defreg{esi}
\defreg{edi}
\defreg{ebp}
\defreg{esp}
\defreg{eip}
\defreg{eflags}

\def\crtrois{\textsc{CR3}}

\newcommand{\ring}[1]{\emph{ring}~#1}

\newcommand{\asminstr}[1]{\texttt{#1}}

\newcommand{\wip}[0]{%
\begin{center}
    \includegraphics[width=5cm]{gen/travaux.pdf}
\end{center}
%\textbf{(wip)}
}

\def\listcons{\texttt{(::)}\xspace}

\tikzstyle{bigbrace} = [decorate,decoration={brace,amplitude=3mm}]

% Debug tikz bounding box
\def\drawbb{\draw[thick] (current bounding box.south west) rectangle (current bounding box.north east);}

% SubFloat + verbatim
\makeatletter
\newbox\sf@box
\newenvironment{SubFloat}[2][]%
{\def\sf@one{#1}%
\def\sf@two{#2}%
\setbox\sf@box\hbox
\bgroup}%
{ \egroup
\ifx\@empty\sf@two\@empty\relax
\def\sf@two{\@empty}
\fi
\ifx\@empty\sf@one\@empty\relax
\subfloat[\sf@two]{\box\sf@box}%
\else
\subfloat[\sf@one][\sf@two]{\box\sf@box}%
\fi}
\makeatother

% comes from the memoir manual {{{

\newcommand*{\addtotextheight}[1]{%
\twocolumn
\addtolength{\textheight}{#1}%
\onecolumn}

%% for main body, bottom of text at 1in, footer below
%% top of header at 1in, first text line double spaced
%% below base of header
\newlength{\linespace}
\setlength{\linespace}{\baselineskip} % current equivalent of \onelineskip
\setlength{\headheight}{\onelineskip}
\setlength{\headsep}{\linespace}
\addtolength{\headsep}{-\topskip}
\setlength{\uppermargin}{1in}
\addtolength{\uppermargin}{\headheight}
\addtolength{\uppermargin}{\headsep}

%% and for the bottom
\setlength{\lowermargin}{1in}
\setlength{\textheight}{\paperheight}

\addtolength{\textheight}{-\uppermargin}
\addtolength{\textheight}{-\lowermargin}
%% footnote settings
\setlength{\footskip}{\onelineskip}
\setlength{\footnotesep}{\onelineskip}


%% the fiddle lengths (..ta.. for title/approval page, others for prelims)
\newlength{\toptafiddle} \setlength{\toptafiddle}{2\linespace}
\newlength{\bottafiddle} \setlength{\bottafiddle}{0pt}
\newlength{\topfiddle}
\setlength{\topfiddle}{\toptafiddle}
\newlength{\botfiddle}
\setlength{\botfiddle}{\onelineskip}

%%% general macro for Abstract, etc., headings
\newcommand*{\pretoctitle}[1]{{\clearpage\centering \vspace*{-\topfiddle}#1\par}}

%%% Start the ACKNOWLEDGEMENTS
\newcommand{\asuacknowledgements}{\pretoctitle{REMERCIEMENTS}}

%%% Start the ABSTRACT
\newcommand{\asuabstract}{ \addtotextheight{-\botfiddle} \pretoctitle{ABSTRACT}}

\addtodef{\mainmatter}{\addtotextheight{\botfiddle}{}}

%% make it easy to center any dedication
\newcommand{\asudedication}[1]{%
{\clearpage\mbox{}\vfill\centering #1 \par\vfill\clearpage}}

% }}} /memoir

% {{{ subtitle command
\newcommand\subtitle[1]{%
  \newcommand\thesubtitle{#1}}
% }}}

% Language name
\def\langname{\textsc{Safespeak}\xspace}

% Center vertically + remove header & footer
\newenvironment{headingpage}{
\thispagestyle{simple}
\vspace*{\fill}
}{
\vspace*{\fill}
}

% Hoare triple
\def\hoare#1#2#3{\ensuremath{\{#1\}~#2~\{#3\}}}

\def\eqdef{\overset{\mathrm{def}}{=\joinrel=}\xspace}

% Lens
\def\mkLens#1#2{\langle #1 | #2 \rangle}
\def\setLens#1#2{\textsc{Lens}_{#1, #2}}
\def\lensGetX#1{\ensuremath{\mathrm{get}_{#1}}}
\def\lensPutX#1{\ensuremath{\mathrm{put}_{#1}}}
\def\lensGet#1#2{\lensGetX{#1}(#2)}
\def\lensPut#1#2#3{\lensPutX{#1}(#2, #3)}
\def\lensNodeBig#1#2{
\vcenter{\hbox{
\begin{tikzpicture}[scale=0.7]
\draw[fill=#1] (-0.5,-0.5) rectangle ++ (1,1);
\draw[fill=#2] (0,0) circle (0.3);
\end{tikzpicture}
}}
}
\def\lensInner#1{
\vcenter{\hbox{
\begin{tikzpicture}[scale=0.7]
\draw[fill=#1] (0, 0) circle (0.3);
\end{tikzpicture}
}}
}
\def\lensNodeBigIdx#1#2#3#4#5{
\vcenter{\hbox{
\begin{tikzpicture}[scale=0.7]
  \draw[fill=#1] (-1.5,-1.5) rectangle ++ (3,3);
  \node[fill=#2, draw, star, minimum size=6mm] at (-0.75, 0.75) {};
  \node[fill=#3, draw, star, minimum size=6mm] at ( 0.75, 0.75) {};
  \node[fill=#4, draw, circle, minimum size=6mm] at (-0.75,-0.75) {};
  \node[fill=#5, draw, regular polygon, regular polygon sides=5, minimum size=6mm] at ( 0.75,-0.75) {};
  \node at (-2,  0.75) {a};
  \node at ( 2,  0.75) {b};
  \node at (-2, -0.75) {c};
  \node at ( 2, -0.75) {d};
\end{tikzpicture}
}}
}
\def\lensInnerStar#1{
\vcenter{\hbox{
\begin{tikzpicture}[scale=0.7]
\node[draw, fill=#1, star] {};
\end{tikzpicture}
}}
}

% Erreurs

\def\serr#1{
  \switch{#1}
    {taint}{Ω_{taint}}
    {array}{Ω_{array}}
    {ptr}{Ω_{ptr}}
    {div}{Ω_{div}}
    {field}{Ω_{field}}
    {var}{Ω_{var}}
    {END}
    % important : étendre la description dans le chapitre 4 et la figure de
    % liste
}

%\def\uGet#1#2{#1 ⇐_U #2}
%\def\uPut#1#2{#1 ⇒_U #2}
\def\cfuF{\textrm{copy\_from\_user}\xspace}
\def\ctuF{\textrm{copy\_to\_user}\xspace}
\def\uGet#1#2{\cfuF(#1, #2)}
\def\uPut#1#2{\ctuF(#1, #2)}

% Litteral notation

\def\eStruct#1{\{#1\}}
\def\eArray#1{[#1]}
\def\tStruct#1{\eStruct{#1}}
\def\eFun#1#2{\mathrm{fun} (#1) \{ #2 \}}

\def\gramorx{\hspace{3mm}\gramor\xspace}
\def\stFun#1{\textsc{Fun}_n}
\def\tComp#1#2{\mathrm{Comp}(#1, #2)}

% Box on top on rules fig
\def\ruleheader#1{
  \begin{minipage}{\linewidth}
    \raggedleft
    \fbox{#1}
  \end{minipage}
}

\def\termFst#1{\ensuremath{\textrm{fst}~#1}}
\def\termSnd#1{\ensuremath{\textrm{snd}~#1}}

\def\ptrtype{\texttt{ptrtype}\xspace}

\def\addrInt{\textrm{addr}^\natural \xspace}
\def\interpPtr#1{\textrm{interp}_{(#1)}}

\tikzset{>=latex}

\def\partccl#1{
  \clearpage~\thispagestyle{cleared}\cleartorecto
  \chapter*{Conclusion de la partie~\ref{#1}}
  \addcontentsline{toc}{chapter}{Conclusion de la partie~\ref{#1}}
}

\def\hexa#1{0x#1}
\def\floatop#1{#1.}

\def\cLookup#1#2{\mathrm{Lookup} (#1, #2)}
\def\cExtend#1#2#3{\mathrm{Extend}(#1, #2 ↦ #3)}
\def\cCleanup#1{\mathrm{Cleanup}(#1)}

\newcommand{\motif}[1]{\cleaders\hbox{#1}\hfil}
\newcommand{\chain}[2]{\leavevmode\hbox to #2{\motif{#1}}}
\newcommand{\diamonds}{$\mkern-.6mu \diamond \mkern-.6mu$}

\def\jolibreak{
\fancybreak{\chain{\diamonds}{0.25\textwidth}}
}

\def\cDom#1{\mathrm{dom}(#1)}

\def\cVarname#1{\mathrm{Varname}(#1)}

\def\cVisible#1{\mathrm{Visible}(#1)}

\def\cRepr#1{\mathrm{Repr}(#1)}

\DeclareUnicodeCharacter{00D7}{\times}        % ×
\DeclareUnicodeCharacter{0393}{\Gamma}        % Γ
\DeclareUnicodeCharacter{03B5}{\epsilon}      % ε
\DeclareUnicodeCharacter{03BB}{\lambda}       % λ
\DeclareUnicodeCharacter{03C3}{\sigma}        % σ
\DeclareUnicodeCharacter{03C4}{\tau}          % τ
\DeclareUnicodeCharacter{2010}{-}             % ‐
\DeclareUnicodeCharacter{2026}{\ldots}        % …
\DeclareUnicodeCharacter{2190}{\leftarrow}    % ←
\DeclareUnicodeCharacter{2191}{\uparrow}      % ↑
\DeclareUnicodeCharacter{2192}{\rightarrow}   % →
\DeclareUnicodeCharacter{21A6}{\mapsto}       % ↦
\DeclareUnicodeCharacter{21D2}{\Rightarrow}   % ⇒
\DeclareUnicodeCharacter{2200}{\forall}       % ∀
\DeclareUnicodeCharacter{2205}{\emptyset}     % ∅
\DeclareUnicodeCharacter{2208}{\in}           % ∈
\DeclareUnicodeCharacter{2260}{\neq}          % ≠
\DeclareUnicodeCharacter{2264}{\le}           % ≤
\DeclareUnicodeCharacter{2265}{\ge}           % ≥
\DeclareUnicodeCharacter{2295}{\oplus}        % ⊕
\DeclareUnicodeCharacter{22A2}{\vdash}        % ⊢
\DeclareUnicodeCharacter{22D8}{\lll}          % ⋘
\DeclareUnicodeCharacter{22D9}{\ggg}          % ⋙

\deflink{cmm}{The C~-~- language}{http://www.cminusminus.org/}
\deflink{ocaml}{OCaml -- Home}{http://ocaml.org/}
\deflink{penjili}{Penjili project}{http://www.penjili.org/}
\deflink{python}{Python Programming Language -- Official Website}{http://www.python.org/}
\deflink{rust}{The Rust Programming Language}{http://www.rust-lang.org/}
\deflink{sparse}{Sparse - a Semantic Parser for {C}}{https://sparse.wiki.kernel.org/index.php/Main_Page}


\usepackage{color,calc,graphicx,soul,fourier}
\definecolor{nicered}{rgb}{.647,.129,.149}
\makeatletter
\newlength\dlf@normtxtw
\setlength\dlf@normtxtw{\textwidth}
\def\myhelvetfont{\def\sfdefault{mdput}}
\newsavebox{\feline@chapter}
\newcommand\feline@chapter@marker[1][4cm]{%
\sbox\feline@chapter{%
\resizebox{!}{#1}{\fboxsep=1pt%
\colorbox{nicered}{\color{white}\bfseries\sffamily\thechapter}%
}}%
\rotatebox{90}{%
\resizebox{%
\heightof{\usebox{\feline@chapter}}+\depthof{\usebox{\feline@chapter}}}%
{!}{\scshape\so\@chapapp}}\quad%
\raisebox{\depthof{\usebox{\feline@chapter}}}{\usebox{\feline@chapter}}%
}
\newcommand\feline@chm[1][4cm]{%
\sbox\feline@chapter{\feline@chapter@marker[#1]}%
\makebox[0pt][l]{% aka \rlap
\makebox[1cm][r]{\usebox\feline@chapter}%
}}
\makechapterstyle{daleif1}{
\renewcommand\chapnamefont{\normalfont\Large\scshape\raggedleft\so}
\renewcommand\chaptitlefont{\normalfont\huge\bfseries\scshape\color{nicered}}
\renewcommand\chapternamenum{}
\renewcommand\printchaptername{}
\renewcommand\printchapternum{\null\hfill\feline@chm[2.5cm]\par}
\renewcommand\afterchapternum{\par\vskip\midchapskip}
\renewcommand\printchaptertitle[1]{\chaptitlefont\raggedleft ##1\par}
}
\makeatother
\chapterstyle{daleif1}

\setlength\beforechapskip{0pt}


\usetikzlibrary{calc}
\usetikzlibrary{shapes}
\usetikzlibrary{arrows}
\usetikzlibrary{decorations.pathreplacing}

\tikzstyle{bigbrace} = [decorate,decoration={brace,amplitude=3mm}]


\usepackage{fancyvrb}

\makeatletter
\def\PY@reset{\let\PY@it=\relax \let\PY@bf=\relax%
    \let\PY@ul=\relax \let\PY@tc=\relax%
    \let\PY@bc=\relax \let\PY@ff=\relax}
\def\PY@tok#1{\csname PY@tok@#1\endcsname}
\def\PY@toks#1+{\ifx\relax#1\empty\else%
    \PY@tok{#1}\expandafter\PY@toks\fi}
\def\PY@do#1{\PY@bc{\PY@tc{\PY@ul{%
    \PY@it{\PY@bf{\PY@ff{#1}}}}}}}
\def\PY#1#2{\PY@reset\PY@toks#1+\relax+\PY@do{#2}}

\expandafter\def\csname PY@tok@gd\endcsname{\def\PY@tc##1{\textcolor[rgb]{0.67,0.00,0.00}{##1}}}
\expandafter\def\csname PY@tok@gu\endcsname{\let\PY@bf=\textbf\def\PY@tc##1{\textcolor[rgb]{0.50,0.00,0.50}{##1}}}
\expandafter\def\csname PY@tok@gt\endcsname{\def\PY@tc##1{\textcolor[rgb]{0.67,0.00,0.00}{##1}}}
\expandafter\def\csname PY@tok@gs\endcsname{\let\PY@bf=\textbf}
\expandafter\def\csname PY@tok@gr\endcsname{\def\PY@tc##1{\textcolor[rgb]{0.67,0.00,0.00}{##1}}}
\expandafter\def\csname PY@tok@cm\endcsname{\let\PY@it=\textit\def\PY@tc##1{\textcolor[rgb]{0.67,0.67,0.67}{##1}}}
\expandafter\def\csname PY@tok@vg\endcsname{\def\PY@tc##1{\textcolor[rgb]{0.67,0.00,0.00}{##1}}}
\expandafter\def\csname PY@tok@m\endcsname{\def\PY@tc##1{\textcolor[rgb]{0.00,0.60,0.60}{##1}}}
\expandafter\def\csname PY@tok@mh\endcsname{\def\PY@tc##1{\textcolor[rgb]{0.00,0.60,0.60}{##1}}}
\expandafter\def\csname PY@tok@cs\endcsname{\let\PY@it=\textit\def\PY@tc##1{\textcolor[rgb]{0.00,0.00,0.67}{##1}}}
\expandafter\def\csname PY@tok@ge\endcsname{\let\PY@it=\textit}
\expandafter\def\csname PY@tok@vc\endcsname{\def\PY@tc##1{\textcolor[rgb]{0.67,0.00,0.00}{##1}}}
\expandafter\def\csname PY@tok@il\endcsname{\def\PY@tc##1{\textcolor[rgb]{0.00,0.60,0.60}{##1}}}
\expandafter\def\csname PY@tok@go\endcsname{\def\PY@tc##1{\textcolor[rgb]{0.53,0.53,0.53}{##1}}}
\expandafter\def\csname PY@tok@cp\endcsname{\def\PY@tc##1{\textcolor[rgb]{0.30,0.51,0.09}{##1}}}
\expandafter\def\csname PY@tok@gi\endcsname{\def\PY@tc##1{\textcolor[rgb]{0.00,0.67,0.00}{##1}}}
\expandafter\def\csname PY@tok@gh\endcsname{\let\PY@bf=\textbf\def\PY@tc##1{\textcolor[rgb]{0.00,0.00,0.50}{##1}}}
\expandafter\def\csname PY@tok@ni\endcsname{\let\PY@bf=\textbf\def\PY@tc##1{\textcolor[rgb]{0.50,0.00,0.00}{##1}}}
\expandafter\def\csname PY@tok@nn\endcsname{\let\PY@ul=\underline\def\PY@tc##1{\textcolor[rgb]{0.00,0.67,0.67}{##1}}}
\expandafter\def\csname PY@tok@no\endcsname{\def\PY@tc##1{\textcolor[rgb]{0.67,0.00,0.00}{##1}}}
\expandafter\def\csname PY@tok@na\endcsname{\def\PY@tc##1{\textcolor[rgb]{0.12,0.56,1.00}{##1}}}
\expandafter\def\csname PY@tok@nb\endcsname{\def\PY@tc##1{\textcolor[rgb]{0.00,0.67,0.67}{##1}}}
\expandafter\def\csname PY@tok@nc\endcsname{\let\PY@ul=\underline\def\PY@tc##1{\textcolor[rgb]{0.00,0.67,0.00}{##1}}}
\expandafter\def\csname PY@tok@nd\endcsname{\def\PY@tc##1{\textcolor[rgb]{0.53,0.53,0.53}{##1}}}
\expandafter\def\csname PY@tok@si\endcsname{\def\PY@tc##1{\textcolor[rgb]{0.67,0.33,0.00}{##1}}}
\expandafter\def\csname PY@tok@nf\endcsname{\def\PY@tc##1{\textcolor[rgb]{0.00,0.67,0.00}{##1}}}
\expandafter\def\csname PY@tok@s2\endcsname{\def\PY@tc##1{\textcolor[rgb]{0.67,0.33,0.00}{##1}}}
\expandafter\def\csname PY@tok@vi\endcsname{\def\PY@tc##1{\textcolor[rgb]{0.67,0.00,0.00}{##1}}}
\expandafter\def\csname PY@tok@nt\endcsname{\let\PY@bf=\textbf\def\PY@tc##1{\textcolor[rgb]{0.12,0.56,1.00}{##1}}}
\expandafter\def\csname PY@tok@nv\endcsname{\def\PY@tc##1{\textcolor[rgb]{0.67,0.00,0.00}{##1}}}
\expandafter\def\csname PY@tok@s1\endcsname{\def\PY@tc##1{\textcolor[rgb]{0.67,0.33,0.00}{##1}}}
\expandafter\def\csname PY@tok@gp\endcsname{\def\PY@tc##1{\textcolor[rgb]{0.33,0.33,0.33}{##1}}}
\expandafter\def\csname PY@tok@sh\endcsname{\def\PY@tc##1{\textcolor[rgb]{0.67,0.33,0.00}{##1}}}
\expandafter\def\csname PY@tok@ow\endcsname{\def\PY@tc##1{\textcolor[rgb]{0.00,0.00,0.67}{##1}}}
\expandafter\def\csname PY@tok@sx\endcsname{\def\PY@tc##1{\textcolor[rgb]{0.67,0.33,0.00}{##1}}}
\expandafter\def\csname PY@tok@bp\endcsname{\def\PY@tc##1{\textcolor[rgb]{0.00,0.67,0.67}{##1}}}
\expandafter\def\csname PY@tok@c1\endcsname{\let\PY@it=\textit\def\PY@tc##1{\textcolor[rgb]{0.67,0.67,0.67}{##1}}}
\expandafter\def\csname PY@tok@kc\endcsname{\def\PY@tc##1{\textcolor[rgb]{0.00,0.00,0.67}{##1}}}
\expandafter\def\csname PY@tok@c\endcsname{\let\PY@it=\textit\def\PY@tc##1{\textcolor[rgb]{0.67,0.67,0.67}{##1}}}
\expandafter\def\csname PY@tok@mf\endcsname{\def\PY@tc##1{\textcolor[rgb]{0.00,0.60,0.60}{##1}}}
\expandafter\def\csname PY@tok@err\endcsname{\def\PY@tc##1{\textcolor[rgb]{0.94,0.00,0.00}{##1}}\def\PY@bc##1{\setlength{\fboxsep}{0pt}\colorbox[rgb]{0.94,0.63,0.63}{\strut ##1}}}
\expandafter\def\csname PY@tok@kd\endcsname{\def\PY@tc##1{\textcolor[rgb]{0.00,0.00,0.67}{##1}}}
\expandafter\def\csname PY@tok@ss\endcsname{\def\PY@tc##1{\textcolor[rgb]{0.00,0.00,0.67}{##1}}}
\expandafter\def\csname PY@tok@sr\endcsname{\def\PY@tc##1{\textcolor[rgb]{0.00,0.60,0.60}{##1}}}
\expandafter\def\csname PY@tok@mo\endcsname{\def\PY@tc##1{\textcolor[rgb]{0.00,0.60,0.60}{##1}}}
\expandafter\def\csname PY@tok@mi\endcsname{\def\PY@tc##1{\textcolor[rgb]{0.00,0.60,0.60}{##1}}}
\expandafter\def\csname PY@tok@kn\endcsname{\def\PY@tc##1{\textcolor[rgb]{0.00,0.00,0.67}{##1}}}
\expandafter\def\csname PY@tok@kr\endcsname{\def\PY@tc##1{\textcolor[rgb]{0.00,0.00,0.67}{##1}}}
\expandafter\def\csname PY@tok@s\endcsname{\def\PY@tc##1{\textcolor[rgb]{0.67,0.33,0.00}{##1}}}
\expandafter\def\csname PY@tok@kp\endcsname{\def\PY@tc##1{\textcolor[rgb]{0.00,0.00,0.67}{##1}}}
\expandafter\def\csname PY@tok@w\endcsname{\def\PY@tc##1{\textcolor[rgb]{0.73,0.73,0.73}{##1}}}
\expandafter\def\csname PY@tok@kt\endcsname{\def\PY@tc##1{\textcolor[rgb]{0.00,0.67,0.67}{##1}}}
\expandafter\def\csname PY@tok@sc\endcsname{\def\PY@tc##1{\textcolor[rgb]{0.67,0.33,0.00}{##1}}}
\expandafter\def\csname PY@tok@sb\endcsname{\def\PY@tc##1{\textcolor[rgb]{0.67,0.33,0.00}{##1}}}
\expandafter\def\csname PY@tok@k\endcsname{\def\PY@tc##1{\textcolor[rgb]{0.00,0.00,0.67}{##1}}}
\expandafter\def\csname PY@tok@se\endcsname{\def\PY@tc##1{\textcolor[rgb]{0.67,0.33,0.00}{##1}}}
\expandafter\def\csname PY@tok@sd\endcsname{\def\PY@tc##1{\textcolor[rgb]{0.67,0.33,0.00}{##1}}}

\def\PYZbs{\char`\\}
\def\PYZus{\char`\_}
\def\PYZob{\char`\{}
\def\PYZcb{\char`\}}
\def\PYZca{\char`\^}
\def\PYZam{\char`\&}
\def\PYZlt{\char`\<}
\def\PYZgt{\char`\>}
\def\PYZsh{\char`\#}
\def\PYZpc{\char`\%}
\def\PYZdl{\char`\$}
\def\PYZti{\char`\~}
% for compatibility with earlier versions
\def\PYZat{@}
\def\PYZlb{[}
\def\PYZrb{]}
\makeatother



\floatstyle{boxed}
\restylefloat{figure}

\setlist[itemize]{label=$\bullet$}

\medievalpage
\checkandfixthelayout

\maxtocdepth{subsection}    % What appears in toc
\maxsecnumdepth{subsection} % What get numbered

%\theoremstyle{definition}
\newtheorem{definition}{Définition}[chapter]
\newtheorem{theorem}{Théorème}[chapter]
\newtheorem{lemma}{Lemme}[chapter]

\def\iffull#1{\if#1\relax\else}
\title{Analyse statique de logiciel système par typage statique fort}
\subtitle{Application au noyau Linux}
\author{Étienne Millon}
\date{xx-yy-2013}

\usepackage{url}
\ifpdf
\usepackage{pdfcolmk}
\fi
%% check if using xelatex rather than pdflatex
\ifxetex
\usepackage{fontspec}
\fi
\usepackage{graphicx}
%%\usepackage{hyperref}
%% drawing package
\usepackage{tikz}
%% for dingbats
\usepackage{pifont}
\providecommand{\HUGE}{\Huge}% if not using memoir
\newlength{\drop}% for my convenience
%% specify the Webomints family
\newcommand*{\wb}[2]{\fontsize{#1}{#2}\usefont{U}{webo}{xl}{n}}
%% select a (FontSite) font by its font family ID
\newcommand*{\FSfont}[1]{\fontencoding{T1}\fontfamily{#1}\selectfont}
%% if you don’t have the FontSite fonts either \renewcommand*{\FSfont}[1]{}
%% or use your own choice of family.
%% select a (TeX Font) font by its font family ID
\newcommand*{\TXfont}[1]{\fontencoding{T1}\fontfamily{#1}\selectfont}
%% Generic publisher’s logo
\newcommand*{\plogo}{\fbox{$\mathcal{PL}$}}

%% Some shades
\definecolor{Dark}{gray}{0.2}
\definecolor{MedDark}{gray}{0.4}
\definecolor{Medium}{gray}{0.6}
\definecolor{Light}{gray}{0.8}
%%%% Additional font series macros
\makeatletter
%%%% light series
%% e.g., kernel doc, section s: line 12 or thereabouts
\DeclareRobustCommand\ltseries
{\not@math@alphabet\ltseries\relax
\fontseries\ltdefault\selectfont}
%% e.g., kernel doc, section t: line 32 or thereabouts
\newcommand{\ltdefault}{l}
%% e.g., kernel doc, section v: line 19 or thereabouts
\DeclareTextFontCommand{\textlt}{\ltseries}
% heavy(bold) series
\DeclareRobustCommand\hbseries
{\not@math@alphabet\hbseries\relax
\fontseries\hbdefault\selectfont}
\newcommand{\hbdefault}{hb}
\DeclareTextFontCommand{\texthb}{\hbseries}
\makeatother


\usetikzlibrary{fadings}

\newcommand*{\titleUL}{
  \newgeometry{left=2cm, right=2cm, top=2cm, bottom=2cm}%
  \begingroup% University of Liege
\drop=0.1\textheight
\begin{center}
{\textsc{Université Pierre et Marie Curie\\
École Doctorale Informatique, Télécommunications et Électronique
}}\\
\vspace{1cm}
\begin{tikzpicture}
  \begin{scope}
    \node (A) {\includegraphics[height=\textwidth]{curies.pdf}};
    \path[opacity=0.2, fill=white]
      (A.north west) rectangle (A.south east);
    %\path[opacity=0.8, fill=white, path fading=north, fading transform={rotate=-45}]
    %(A.north west) rectangle (A.south east);
    %\clip (A.center) circle(8.5cm);
    %\shade[opacity=0.2, shading=axis,bottom color=black!10,top color=black,shading angle=-45]
    %(A.south west) rectangle (A.north east);
  \end{scope}
\end{tikzpicture}
\vspace{-13cm}
\\[\drop]
\begin{minipage}{13cm}
  \begin{center}
\rule{\textwidth}{1pt}\par
\vspace{0.5\baselineskip}
{\huge\bfseries \thetitle\\
\vspace{5mm}
\large --- \thesubtitle ---
}\\[0.5\baselineskip]
\rule{\textwidth}{1pt}
  \end{center}
\end{minipage}
\par
\vspace{1cm}
{{\Large\textsc{\theauthor}} \\
sous la direction d'Emmanuel Chailloux et de Sarah Zennou
}
\vspace{16mm}\\
{{\LARGE\textsc{Thèse}}\\
pour obtenir le titre de\\
{\Large Docteur en Sciences}\\
{\large mention Informatique}
}
\vspace{5mm}\\
{
  Version du \today{} \\
}
{\large Soutenue le \thedate{} devant un jury composé de\\

\textbf{Président}\\
---\\
\textbf{Rapporteurs}\\
---\\
---\\
%Sandrine Blazy\\
%Pierre Jouvelot\\
\textbf{Directeurs}\\
Emmanuel Chailloux\\
Sarah Zennou\\
\textbf{Examinateurs}\\
---\\
---
}
\end{center}
\endgroup
\restoregeometry
}


\begin{document}

\frontmatter

\begin{titlingpage}
\titleUL{}
\end{titlingpage}

\cleartorecto

\asuabstract

\cleartoverso
\newgeometry{margin=2cm}

\begin{headingpage}
\thispagestyle{empty}

\begin{center}\Large \textbf{Résumé}\end{center}

Les noyaux de système d'exploitation manipulent des données fournies par les
programmes utilisateur via certaines interfaces, comme les appels système. Si
celles-ci sont manipulées sans prendre une attention particulière, une faille de
sécurité connue sous le nom de \emph{Confused Deputy Problem} peut amener à des
fuites de données confidentielles ou l'elévation de privilège d'un attaquant.

Le but de cette thèse est d'utiliser des techniques de typage statique afin de
détecter les manipulations dangereuses de pointeurs contrôlés par l'espace
utilisateur.

De nombreux systèmes d'exploitation, dont le noyau Linux sur lequel nous nous
basons, sont écrits dans le langage C. Afin de pouvoir appliquer les techniques
de typage statique à ce langage, on commence par isoler un sous-langage sûr
nommé \langname, pour lequel on définit une sémantique opérationnelle et un
système de types. Les propriétés classiques de sûreté du typage sont établies.

Dans un second temps on étend \langname pour intégrer la notion de valeur
provenant de l'espace utilisateur. On montre que l'extension triviale du système
de types brise la sûreté du typage, et que l'on peut réétablir celle-ci en
donnant un type particulier aux pointeurs dont la valeur est contrôlée par
l'espace utilisateur, forçant leur déférencement à se faire au sein d'un
ensemble limité de fonctions sûres. Un cas d'étude est déroulé pour montrer que
cette technique permet de détecter un bug qui a frappé un pilote de carte
graphique AMD dans le noyau Linux.

Ces travaux ouvrent des perspectives sur deux points. D'une part, il est
possible d'appliquer cette analyse à de plus grandes parties du noyau Linux. Les
diverses interfaces qui exposent des données contrôlées par l'espace utilisateur
peuvent bénéficier de cette vérification statique. Ceci est d'autant plus vrai
que les pilotes de périphériques dont elles sont issues sont en général moins
testées que le reste du code du noyau. D'autre part, cette technique revient à
introduire des types abstraits opaques, séparant leur représentation de leur
interface. Ce concept étant absent de C, de nombreuses erreurs de programmation
notamment liées aux manipulations d'entiers peuvent également être vérifiées à
l'aide de cette méthode.

\begin{center}\Large \textbf{Abstract}\end{center}

Lorem ipsum dolor sit amet, consectetur adipiscing elit. Donec gravida, diam
ullamcorper volutpat varius, est dui blandit nibh, et posuere nisi tortor ut
lectus. Aliquam in ipsum at nisl luctus pulvinar. Nulla pulvinar accumsan
interdum. Aliquam varius est et ligula commodo varius quis gravida justo. Ut
sit amet eros orci, at iaculis metus. Cras placerat condimentum magna ut
dictum. Mauris id sollicitudin quam. Ut pretium aliquet lacus, et elementum
odio blandit quis.

Class aptent taciti sociosqu ad litora torquent per conubia nostra, per
inceptos himenaeos. Morbi tincidunt, justo sed venenatis imperdiet, diam dui
blandit nisl, et sodales urna mauris sit amet felis. Sed congue sodales
sollicitudin. Integer imperdiet tellus et ligula consequat quis commodo metus
posuere. Donec id vehicula mi. Nullam id tincidunt orci. Nunc luctus ligula ut
urna laoreet consequat. In hac habitasse platea dictumst. Sed imperdiet
ultricies diam id cursus. Quisque metus orci, tempor et dapibus eget,
ullamcorper in augue. Nullam non urna metus, id accumsan lacus. Duis erat
arcu, porttitor non laoreet vitae, ultrices vitae tortor. Pellentesque
interdum faucibus tellus eu semper. Nulla ac aliquet mauris.

In in metus purus. Nunc ullamcorper posuere turpis, vel laoreet tortor posuere
at. Ut euismod arcu sit amet dui varius vel varius neque pellentesque. Lorem
ipsum dolor sit amet, consectetur adipiscing elit. Aliquam nec mi interdum
ipsum eleifend tincidunt. Duis bibendum sem vitae risus facilisis mattis. Sed
varius eleifend nulla vel commodo. Morbi risus metus, egestas non commodo id,
ullamcorper a enim. Integer et augue et orci feugiat rhoncus a in augue.

Nulla elementum felis eu eros pretium quis porttitor felis faucibus. Vivamus a
vulputate tellus. Suspendisse potenti. Morbi dapibus, quam ac consectetur
hendrerit, lorem urna faucibus dui, in vestibulum tellus nisl sit amet ligula.
Morbi ultricies venenatis nisi sit amet scelerisque. Nam pharetra bibendum
lacinia. Suspendisse potenti. Suspendisse elementum tellus et mi hendrerit
sagittis.

Ut dignissim fermentum cursus. Praesent feugiat venenatis rutrum. Cras eget
nibh id quam ultrices pellentesque eget non elit. Maecenas justo lorem,
fermentum vitae dictum eu, egestas eu urna. In hac habitasse platea dictumst.
Cras nec arcu at purus aliquet accumsan. Nam interdum malesuada orci et
sagittis

% TODO (PLAS suit)

%Manipulating user-provided pointers in the kernel of an operating system can
%lead to security flaws if done in an incautious manner. We present an efficient
%system to detect and prevent this class of erroneous memory manipulation.

%At the core of our approach is \langname, an imperative language that we equip
%with a qualified type system, where two kinds of pointers are distinguished:
%\emph{safe} pointers, whose value is statically proved to be controlled by the
%kernel, and \emph{unsafe} ones, whose value comes from userspace through
%run-time system calls. Dereferencing unsafe pointers is forbidden in a static
%manner by the means of a strong type system.

%A concrete case study is described based on a bug that affected a video driver
%in the Linux kernel. We also explain a technique to automatically translate GNU
%C code to our core language, which will enable us to analyze larger fractions of
%the kernel in order to find similar vulnerabilities.

\end{headingpage}
\restoregeometry


\asuacknowledgements

Remerciements


\asudedication{Cassedédi}

\tableofcontents

\mainmatter

\chapter{Introduction}
Communication,
audiovisuel,
transports,
médecine:
tout ces domaines se sont transformés dans les dernières décennies,
en particulier grâce à la révolution numérique.
En effet le plus petit appareil électrique contient maintenant des composants
matériels programmables.

% TODO faire une intro "la vie est numérique" avec du Berry @ CdF

En 2013, on pense bien sûr aux téléphones portables dont la fonctionnalité et la
complexité les rapprochent des ordinateurs de bureau. Par exemple, le système
d'exploitation Android de Google est basé sur le noyau Linux, destiné à la base
aux micro-ordinateurs.

Le noyau d'un système d'exploitation est chargé de faire l'intermédiaire entre
le matériel (processeur, mémoire, périphériques, …) et les applications
exécutées sur celui-ci (par exemple un navigateur web, une calculatrice ou un
carnet d'adresses). Il doit aussi garantir la sécurité de celles-ci.

En tant qu'intermédiaire de confiance, le noyau a un certain nombre de
responsabilités et est le seul à avoir accès à certaines informations sensibles.
Il est capital de s'assurer qu'il est bien le seul à pouvoir y accéder. En
particulier, il faut pouvoir vérifier que les requêtes faites par l'utilisateur
au noyau ne peuvent pas volontairement ou involontairement détourner ce dernier
et lui faire fuiter des informations confidentielles.

Le problème est que comme tous les logiciels, les noyaux de système
d'exploitation sont écrits par des humains, qui ne sont pas parfaits. Loin de
là: on estime qu'avant relecture, 1000 lignes de code contiennent entre 5 et 100
erreurs de programmation en moyenne.
% TODO sourcer
Les activités de relecture et de
débogage ont beau prendre la majeure partie du temps de développement, il est
facile de laisser passer des défauts de programmation.

Une technique efficace est de réaliser des tests, c'est-à-dire exécuter le
programme sous un environnement contrôlé. On peut alors détecter des
comportements non désirés. Mais même avec des tests les plus exhaustifs il n'est
pas possible de couvrir tous les cas d'utilisation.

Une autre approche est d'analyser le code source du programme avant de
l'exécuter et de refuser de lancer les programmes qui contiennent des erreurs,
quitte à limiter la possibilité d'en écrire certains. C'est l'analyse statique
de programmes.

Une des techniques d'analyse statique les plus répandues est le typage statique,
qui consiste à associer à chaque morceau de programme, une étiquette
décrivant quel genre de valeur sera produite par son évaluation. Par exemple, si
$n$ est le nom d'une variable entière, alors $n + 2$ produira toujours une
valeur entière.

Pour garantir l'isolation d'un noyau de système d'exploitation, un des points
cruciaux est de restreindre la manière dont sont traitées les informations
provenant des programmes utilisateur. Le but de cet thèse est de montrer que le
typage statique peut être utilisé pour détecter et interdire ces manipulations
dangereuses.

\section{Rôle d'un système d'exploitation}

Un ordinateur est constitué de nombreux composants matériels: microprocesseur,
mémoire, et divers périphériques. Et au niveau de l'utilisateur, des dizaines de
logiciels permettent d'effectuer toutes sortes de calculs et de communications.
Le système d'exploitation permet de faire l'interface entre ces niveaux
d'abstraction.

Au cours de l'histoire des systèmes informatiques, la manière de les programmer
a beaucoup évolué. Au départ, les programmeurs avaient accès au matériel dans
son intégralité: toute la mémoire pouvait être accédée, toutes les instructions
pouvaient être utilisées.

Néanmoins c'est un peu restrictif, puisque cela ne permet qu'à une personne
d'interagir avec le système. Dans la seconde moitié des années 60, sont apparus
les premiers systèmes ``à temps partagé'', permettant à plusieurs utilisateurs
de travailler en même temps.

Permettre l'exécution de plusieurs programmes en même temps est une idée
révolutionnaire, mais elle n'est pas sans difficultés techniques: en effet les
ressources de la machine doivent être aussi partagées entre les utilisateurs et
les programmes. Par exemple, plusieurs programmes vont utiliser le processeur
les uns à la suite des autres; et chaque programme aura à sa disposition une
partie de la mémoire principale, ou du disque dur.

Si plusieurs programmes s'exécutent de manière concurrente sur le même matériel,
il faut s'assurer que l'un ne puisse pas écrire dans la mémoire de l'autre, ou
que les deux n'utilisent pas la carte réseau en même temps. Ce sont des rôles du
système d'exploitation.

Ainsi, au lieu d'accéder directement au matériel via des instructions de bas
niveau, les programmes communiquent avec le noyau, qui centralise donc les
appels au matériel, et abstrait certaines opérations.

Par exemple, comparons ce qui se passe concrètement lors de la copie de données
depuis un cédérom ou une clef USB.\@

\begin{itemize}

  \item Dans le cas du cédérom, il faut interroger le bus SATA, interroger le
    lecteur sur la présence d'un disque dans le lecteur, activer le moteur,
    calculer le numéro de trame des données sur le disque, demander la lecture,
    puis déclencher une copie de la mémoire.

  \item Avec une clef, il faut interroger le bus USB, rechercher le bon numéro
    de périphérique, le bon numéro de canal dans celui-ci, lui appliquer une
    commande de lecture au bon numéro de bloc, puis copier la mémoire.

\end{itemize}

Ces deux opérations, bien qu'elles aient la même intention (copier de la mémoire
depuis un périphérique amovible), ne sont pas effectuées en extension de la même
manière. C'est pourquoi le système d'exploitation fournit les notions de
fichier, lecteur, etc: le programmeur n'a plus qu'à utiliser des commandes de
haut niveau (``monter un lecteur'', ``ouvrir un fichier'', ``lire dans un
fichier'') et selon le type de lecteur, le système d'exploitation effectuera les
actions appropriées.

En résumé, un système d'exploitation est l'intermédiaire entre le logiciel et le
matériel, et en particulier est responsable de la gestion de la mémoire, des
périphériques et des processus. Les détails d'implantation ne sont pas présentés
à l'utilisateur; à la place il manipule des abstractions, comme la notion de
fichier. Pour une explication détaillée du concept de système d'exploitation
ainsi que des cas d'étude, on pourra se référer à~\cite{tanenbaum}.

\section{Séparation entre noyau et espace utilisateur}

Puisque le noyau est garant du bon fonctionnement du système, il ne doit pas
pouvoir être manipulé directement par l'utilisateur ou les programmes exécutés.
Ainsi, il est nécessaire de mettre en place des protections entre les espaces
noyau et utilisateur.

Au niveau matériel, on utilise la notion de \emph{niveaux de privilèges} pour
déterminer s'il est possible d'exécuter une instruction.

D'une part, le processeur contient un niveau de privilège intrinsèque. D'autre
part, chaque zone mémoire contenant du code ou des données possède également un
niveau de privilège minimum nécessaire. L'exécution d'une instruction est alors
possible si et seulement si le niveau de privilège du processeur est supérieur à
celui de l'instruction et des opérandes mémoires qui y sont présentes
\footnote{
  Ici ``supérieur'' est synonyme de ``plus privilégié''. Dans l'implantation
  d'Intel présentée dans le chapitre~\ref{cha:os}, les niveaux sont numérotés de
  0 à 3 où le niveau 0 est le plus privilégié.
}.

Par exemple, supposons qu'un programme utilisateur contienne l'instruction
\texttt{mov \%eax, \hexa{54c3b2f7}}. Celle-ci consiste à déplacer le contenu du
registre \eax{} vers l'adresse mémoire \texttt{\hexa{54c3b2f7}}. Puisque cette
adresse est dans l'espace utilisateur (la raison est décrite dans le
chapitre~\ref{cha:os}), aucune erreur de protection mémoire n'est déclenchée.

Ainsi, pour une instruction manipulant des données en mémoire, les accès
possibles sont décrits dans la figure~\ref{fig:erreursec}. En cas
d'impossibilité (signalée par un \Square), une erreur se produit et l'exécution
s'arrête.

\begin{figure}
\begin{center}
\def\modeK{Noyau\xspace}
\def\modeU{Utilisateur\xspace}
\begin{tabular}{cccc}
  Mode du processeur
& Privilège (code)
& Privilège (données)
& Accès possible \\
\midrule
  \modeK{} & \modeK{} & \modeK{} & \CheckedBox{} \\
  \modeK{} & \modeK{} & \modeU{} & \CheckedBox{} \\
  \modeK{} & \modeU{} & \modeK{} & \CheckedBox{} \\
  \modeK{} & \modeU{} & \modeU{} & \CheckedBox{} \\
  \modeU{} & \modeK{} & \modeK{} & \Square{} \\
  \modeU{} & \modeK{} & \modeU{} & \Square{} \\
  \modeU{} & \modeU{} & \modeK{} & \Square{} \\
  \modeU{} & \modeU{} & \modeU{} & \CheckedBox{} \\
\end{tabular}
\end{center}

\caption{ Règles d'exécution de code et d'accès à la mémoire. }
\label{fig:erreursec}

\end{figure}

En plus de cette vérification, certains types d'instructions sont explicitement
réservés au mode le plus privilégié: par exemple les lectures ou écritures sur
des ports matériels, ou celles qui permettent de définir les niveaux de
privilèges des différentes zones mémoire.

Les programmes utilisateur ne pouvant pas accéder à ces instructions de bas
niveau, ils sont très limités dans ce qu'ils peuvent faire. Puisque
l'interaction avec le matériel (comme l'écriture sur un disque dur) se fait
uniquement via des instructions privilégiées, ils sont limités à l'utilisation
du processeur et de la mémoire, permettant uniquement de réaliser des calculs.

Pour utiliser le matériel ou accéder à des abstractions de haut niveau (comme
créer un nouveau processus), ils doivent donc passer par l'intermédiaire du
noyau. La communication entre le noyau et les programmes utilisateur est
constituée par le mécanisme des \emph{appels système}.

Lors d'un appel système, une fonction du noyau est invoquée (en mode noyau) avec
des paramètres provenant de l'utilisateur. Il faut donc être particulièrement
précautionneux dans le traitement de ces données.

Par exemple, considérons un appel système de lecture depuis un disque: on passe
au noyau les arguments $(d, o, n, a)$ où $d$ est le nom du disque, $o$ (pour
\emph{offset}) l'adresse sur le disque où commencer la lecture, $n$ le nombre
d'octets à lire et $a$ l'adresse en mémoire où commencer à stocker les
résultats.

Dans le cas d'utilisation prévu, le noyau va copier la mémoire lue dans $a$. Le
processeur est en mode noyau, en train d'exécuter une instruction du noyau
manipulant des données utilisateur. D'après la figure~\ref{fig:erreursec} aucune
erreur ne se produit.

Mais même si ce cas ne produit pas d'erreur à l'exécution, il est tout de même
erroné. En effet, si on passe à l'appel système une adresse $a$ faisant partie
de l'espace noyau, que se passe-t-il?

L'exécution est presque identique: au moment de la copie on est en mode noyau,
en train d'exécuter une instruction du noyau manipulant des données noyau.
Encore une fois il n'y a pas d'erreur à l'exécution.

On peut donc écrire n'importe où en mémoire. De même, une fonction d'écriture
sur un disque (et lisant en mémoire) permettrait de lire de la mémoire du noyau.
À partir de ces primitives, on peut accéder aux autres processus exécutés, ou
détourner l'exécution vers du code arbitraire. L'isolation est totalement
brisée à cause de ces appels système.

La cause de ceci est qu'on a accédé à la mémoire en testant les privilèges du
noyau au lieu de tester les privilèges de celui qui a fait la requête
(l'utilisateur). Ce problème est connu sous le nom de \emph{confused deputy
problem}~\cite{hardy88confused}.

La bonne manière d'implanter un appel système est donc d'interdire le
déréférencement direct des pointeurs dont la valeur peut être contrôlée par
l'utilisateur. Dans le cas du passage par adresse d'un argument, il aurait fallu
vérifier à l'exécution que celui-ci a bien les mêmes privilèges que l'appelant.

Il est facile d'oublier d'ajouter cette vérification, puisque le cas ``normal''
fonctionne. Avec ce genre d'exemple on voit comment les bugs peuvent arriver si
fréquemment et pourquoi il est aussi capital de les détecter avant l'exécution.

\section{Systèmes de types}
\label{sec:types-intro}

La plupart des langages de programmation incorporent la notion de type, dont un
des buts est d'empêcher de manipuler des données incompatibles entre elles.

Au niveau du langage machine, les seules données qu'un ordinateur manipule sont
des nombres. Selon les opérations effectuées, ils seront interprétés comme des
entiers, des adresses mémoires, ou des caractères. Pourtant il est clair que
certaines opérations n'ont pas de sens: par exemple, multiplier un nombre par
une adresse, ou déréférencer le résultat d'une division sont des comportements
qu'on voudrait pouvoir empêcher.

En un mot, le but du typage est de classifier les objets et de restreindre les
opérations possibles selon la classe d'un objet: en somme, ``ne pas ajouter des
pommes et des oranges''. Le modèle qui permet cette classification est appelé
\emph{système de types} et est en général constitué d'un ensemble de
\emph{règles de typage}, comme ``un entier plus un entier égale un entier''.

\paragraph{Typage dynamique}

Dans ce cas, chaque valeur manipulée par le programme est décorée d'une
étiquette définissant comment interpréter la valeur en question. Les règles de
typage sont alors réalisées à l'exécution. Par exemple, l'opérateur ``$+$''
vérifie que ces deux opérandes ont une étiquette ``entier'', et construit alors
une valeur obtenue en faisant l'addition des deux valeurs, avec une étiquette
``entier''. Par exemple, le langage Python~\link{python} utilise cette
stratégie.

\paragraph{Typage statique}

Dans ce cas on fait les vérifications à la compilation. En quelque sorte,
l'approche dynamique est pessimiste, puisqu'elle demande de traiter à chaque
étape le cas où les types ne sont pas corrects. Intuitivement, dans le cas où
toutes les fonctions se comportent bien, faire la vérification est inutile. Pour
vérifier ceci, on donne à chaque fonction un contrat comme ``si deux entiers
sont passés, et que la fonction renvoie une valeur, alors cette valeur sera un
entier''. Cet ensemble de contrats peut être vérifié statiquement par le
compilateur, à l'aide d'un système de types statique.

Par exemple, on peut dire que l Ainsi l'opérateur ``$+$'' a pour type $(\tInt,
\tInt) → \tInt$. Cela veut dire que si on lui passe deux entiers ($\tInt$
\tInt), alors la valeur obtenue est également un entier.

\paragraph{Typage fort ou faible}

Contrairement à la distinction claire entre typage statique ou dynamique, la
séparation entre typage fort et faible est moins nette. En poussant à l'extrême,
les systèmes de types forts permettent d'éliminer totalement la nécessité de
réaliser des tests de typage. Mais souvent ce n'est pas le cas, car il peut y
avoir des constructions au sein du langage qui permettent de contourner le
système de types, comme un opérateur de transtypage.

\paragraph{Polymorphisme}

Parfois, il est trop restrictif de donner un unique type à une fonction. Si on
considère une fonction ajoutant un élément à une liste, ou une autre
triant un tableau en place, leur type doit-il faire intervenir le type des
éléments manipulés?

En première approximation, on peut imaginer fournir une version du code par type
de données à manipuler. C'est la solution retenue par les premières versions du
langage Pascal, ce qui rendait très difficile l'écriture de
bibliothèques~\cite{PascalNoEscape}. On parle alors de monomorphisme.

Une autre manière de procéder est d'autoriser plusieurs fonctions à avoir le
même nom, mais avec des types d'arguments différents. Par exemple, on peut
définir séparément l'addition entre deux entiers, entre deux flottants, ou entre
un entier et un flottant. Selon les informations connues à la compilation, la
bonne version sera choisie. C'est ainsi que fonctionnent les opérateurs en C++.
On parle de polymorphisme \emph{ad hoc}, ou de surcharge.

La dernière possibilité est le polymorphisme paramétrique, qui consiste à
utiliser le même code quelque soit le type des arguments. Dans ce cas, on
utilise une seule fonction pour traiter une liste d'entiers ou une liste de
flottants, par exemple. Au lieu d'associer à chaque fonction un type, dans
certains cas on lui associe un type paramétré, instanciable en un type concret.
Dans le cas des fonctions de traitement de liste, l'idée est que lorsqu'on ne
touche pas aux éléments, alors le traitement est valable quelque soit leur type.
Cette technique a été décrite en premier dans~\cite{Milner78}.

Pour un tour d'horizon de différents systèmes types de statiques, avec en
particulier du polymorphisme, on pourra se référer à~\cite{TAPL}.

\section{Langages d'implantation}

Puisque notre but est d'analyser du code provenant de systèmes d'exploitation,
il est nécessaire de s'intéresser aux langages de programmation dans lesquels
ils sont écrits.

\paragraph{Assembleur}

Le noyau d'un système d'exploitation nécessite d'accéder au matériel, donc il
est naturellement bas niveau. Pour accéder aux fonctionnalités spécifiques de
chaque processeur, il est donc nécessaire d'en implanter une partie dans le
langage d'assemblage natif de chaque architecture. Historiquement, les premiers
système d'exploitations étaient entièrement écrits en assembleur. Cela est de
plus en plus rare, sauf dans les cas où les ressources sont trop limitées pour
exécuter du code compilé, comme dans les systèmes embarqués ou temps-réel.

\paragraph{C}

Le système Unix, développé à partir de 1969, a tout d'abord été développé en
assembleur sur un mini-ordinateur PDP-7, puis a été porté sur d'autres
architectures matérielles. Pour aider ce portage, il a été nécessaire de créer
un ``assembleur portable'', le langage C~\cite{KandR,AnsiC}. Son but est de
fournir des abstractions au dessus du langage d'assemblage. Les structures de
contrôle (\texttt{if}, \texttt{while}, \texttt{for}) permettent d'utiliser la
programmation structurée, c'est-à-dire en limitant l'utilisation de
l'instruction \texttt{goto}. Les types de données sont également abstraits de la
machine: ainsi, \texttt{int} désigne un entier machine, indépendamment de sa
taille concrète. Son système de types est assez rudimentaire: toutes les formes
de transtypage sont acceptées, certaines conversions sont insérées
automatiquement par le compilateur, et la plupart des abstractions fournies par
le langage sont perméables. Le noyau Linux est écrit dans un dialecte du langage
C. Le noyau du système Mac OSX d'Apple est également un dérivé d'Unix, et est
donc aussi écrit dans ce langage.

\paragraph{C++}

Le manque de modularité de C a conduit à la création d'un nouveau langage plus
riche autour de celui-ci. Dans les années 1980, C++ est apparu comme un C
amélioré, avec de nombreuses fonctionnalités supplémentaires comme un modèle
objet avec des classes, du polymorphisme \emph{ad hoc} et un système de
\emph{templates} puissant permettant d'implanter une forme de polymorphisme
universel. Microsoft Windows, initialement écrit en C, s'est enrichi de plus en
plus de code C++, en faisant un des principaux systèmes d'exploitation écrits
dans ce langage.

\paragraph{Ada}

Le langage C est bas niveau, mais son système de types n'est pas suffisamment
évolué pour le rendre plus sûr que le langage d'assemblage vers lequel il est
compilé. Dans les années 1970, le département de la défense américain
(\emph{DoD}) a créé un comité dans le but de créer un nouveau langage pour ses
systèmes temps-réels embarqués. Le résultat est Ada, un langage impératif avec
un système de types statique plus contraignant mais plus sûr que celui de C. Par
exemple, deux types entiers différents (comme les entiers de 16 bits et ceux de
32 bits) ne sont pas compatibles. Il faut ajouter une conversion explicite dans
le code, qui se traduira par un test à l'exécution et éventuellement par une
erreur \texttt{Constraint\_Error}. Bien qu'il ne se soit pas beaucoup développé
dans les applications grand public, Ada garde toujours une place de choix dans
l'implantation de systèmes embarqués, notamment dans les applications spatiales
et de défense.

\paragraph{Divers} % TODO titre

Il est bien sûr possible d'implanter un système d'exploitation dans n'importe
quel langage, si celui-ci peut être compilé vers du langage machine et qu'il
permet facilement d'inclure des fragments d'assembleur pour les parties de bas
niveau. Néanmoins il s'agit en général de travaux de recherche. On peut ainsi
trouver des systèmes d'exploitation écrits en C\#~\cite{SingularityOverview},
Haskell~\cite{HaskellOS} ou Ocaml~\cite{unikernels}.

% TODO[B] TROLL

\section{Langages d'analyse}

Les langages décrits précédemments sont faits pour être facilement écrits par
des programmeurs humains. En général ils sont ambigüs ou ont des comportements
implicites. Pour analyser du code source, il est plus pratique d'avoir une
représentation intermédiaire plus simple afin d'avoir moins de traitements
dupliqués. Dans de nombreux projets, des sous-ensembles de C ont été définis
pour aller dans ce sens.

\subsection*{Middle-ends}

Les premiers candidats sont bien entendu les représentations intermédiaires
utilisées dans les compilateurs C. Elles ont l'avantage d'accepter en plus du C
standard, les diverses extensions (GNU, Microsoft, Plan9) utilisées par la
plupart des logiciels. En particulier, le noyau Linux repose fortement sur les
extensions GNU.\@

\paragraph{GCC} utilise une représentation interne nommée
GIMPLE\cite{gcc-gimple}. Il s'agit d'une structure d'arbre écrite en C, reposant
sur de nombreuses macros afin de cacher les détails d'implantation interne
pouvant varier entre deux versions. Cette représentation étant réputée difficile
à manipuler, le projet MELT\cite{gcc-melt} permet de générer un greffon de
compilateur à partir d'un dialecte de Lisp.

\paragraph{LLVM}\cite{llvm-pres} est un compilateur développé par la communauté
puis sponsorisé Apple. À la différence de GCC, sa base de code est écrite en
C++. Il utilise une représentation intermédiaire qui peut être manipulée soit
sous forme d'une structure de données C++, soit d'un fichier de code-octet
compact, soit sous forme textuelle.

\paragraph{Cmm} est une représentation interne utilisée pour la génération de
code lors de la compilation d'Objective Caml~\link{ocamlManual}, et disponible
dans les sources du compilateur sous le chemin \texttt{asmcomp/cmm.mli} (il
s'agit donc d'une structure de données OCaml). Ce langage a l'avantage d'être
très restreint, mais malheureusement il n'existe pas directement de traducteur
permettant de compiler C vers Cmm.

\paragraph{C-~-}\cite{spjcmm} \link{cmm}, dont le nom est inspiré du précédent,
est un projet qui visait à unifier les langages intermédiaires utilisés par les
compilateurs. L'idée est que si un front-end peut émettre du C-~- (sous forme de
texte), il est possible d'obtenir du code machine efficace. Le compilateur
Haskell GHC, par exemple, utilise une représentation intermédiaire très
similaire à C-~-.

\subsection*{Langages intermédiaires ad hoc}

Comme le problème de construire une représentation intermédiaire adaptée à une
analyse statique n'est pas nouveau, plusieurs projets ont déjà essayé d'y
apporter une solution. Puisque qu'ils sont développés en parallèle des
compilateurs, le support des extensions est en général moins important dans ces
langages.

\paragraph{CIL}\cite{NeculaCil} \link{kerneisCil} est une représentation en
OCaml d'un programme C, développée depuis 2002. Grâce à un mécanisme de
greffons, elle permet de prototyper rapidement des analyses statiques de
programmes.

\paragraph{Compcert} est un projet qui vise à produire un compilateur certifié
pour C. C'est-à-dire que le fait que les transformations conservent la
sémantique est prouvé. Il utilise de nombreux langages intermédiaires, dont
CIL.\@ Pour le front-end, le langage se nomme Clight\cite{cfront}. Les passes de
middle-end, quant à elles, utilisent Cminor\cite{cminorSL}.

\paragraph{Newspeak}\cite{newspeak} est un langage intermédiaire développé par
EADS Innovation Works, et qui est spécialisé dans l'analyse de valeurs par
interprétation abstraite. Il sera décrit plus en détails dans la
Section~\ref{sec:npk}.

\subsection*{La solution retenue: \langname{}}

Dans la suite de cette thèse nous nous appuierons sur \langname{}, un langage
intermédiaire minimal, pour exprimer les règles de typage de la manière la plus
simple possible.

\section{Plan de la thèse}

Cette thèse comporte trois parties.

\paragraph{La partie~\ref{part:ctx}} présente le contexte de ces travaux. Le
fonctionnement général d'un système d'exploitation y est détaillé, et les
problèmes de manipulation de pointeurs contrôlés par l'utilisateur y sont
introduits. On fait ensuite un tour d'horizon des techniques existantes
permettant de traiter ce problème. Cette partie se conclut par la proposition au
cœur de cette thèse.

\paragraph{La partie~\ref{part:lang}} décrit notre solution: \langname{}, un
langage impératif. On y décrit sa syntaxe, sa sémantique ainsi qu'un système de
types statiques. On l'étend ensuite pour capturer les problèmes d'adressage
mémoire présents dans les systèmes d'exploitation en ajoutant des pointeurs
contrôlés par l'utilisateur. Le système de types est également étendu. Pour
chacune de ces variantes, on établit la propriété de sûreté de typage reliant la
sémantique d'exécution aux types statiques.

\paragraph{La partie~\ref{part:xp}} documente la démarche expérimentale associée
à ces travaux. L'implantation du système de types est décrite, afin que la
manière de transformer automatiquement du code C en \langname{}. Un cas d'étude
est déroulé, consistant d'un bug ayant touché le noyau Linux. Il est démontré
que le système de type capture précisément ce genre d'erreur de programmation.
Enfin, les possibilités d'extension tant théoriques qu'expérimentales sont
présentées.

%  - historique des techniques
%  - analyse non intrusive
%  - méthodes de développement du noyau
%  - syntaxe, sémantique, typage

% TODO [E] + besoins EADS

% vim: spelllang=fr


\part{Méthodes formelles pour la sécurité}

\chapter{Systèmes d'exploitation}

\label{cha:os}
\section{Rôle d'un système d'exploitation}

% Fig Archi {{{
\begin{figure}
\centering
\begin{tikzpicture}
[ every node/.style={}
, node distance=2cm
]

\node (cpu) {CPU};

\node[below of=cpu] (northb) {Northbridge};

\node[right of=northb] (mem) {Mémoire};
\node[left of=northb]  (agp) {AGP};
\node[below of=northb] (southb) {Southbridge};

\node[below left of=southb, xshift=-1cm] (sata) {SATA};
\node[below right of=southb, xshift=1cm] (usb)  {USB};

\node[below left of=sata]  (hdd)     {Disque dur};
\node[below right of=sata] (cdrom)   {CD-ROM};

\node[below left of=usb]  (memo)    {Mémorette};
\node[below right of=usb] (clavier) {Clavier};

\draw (cpu) -- (northb);
\draw (northb) -- (agp);
\draw (northb) -- (mem);
\draw (southb) -- (northb);
\draw (southb) -- (sata);
\draw (southb) -- (usb);
\draw (usb) -- (memo);
\draw (usb) -- (clavier);
\draw (sata) -- (hdd);
\draw (sata) -- (cdrom);

\end{tikzpicture}

\caption{Architecture simplifiée d'un ordinateur}
\label{fig:archi-simplifiee}
\end{figure}
% }}}

Au plus bas, un ordinateur est constitué de nombreux composants matériels :
micro-processeur, mémoire, et divers périphériques
(figure~\ref{fig:archi-simplifiee}). Et au niveau de l'utilisateur, de dizaines
de logiciels permettant d'effectuer toutes sortes de calculs et de
communication. Le système d'exploitation permet de faire l'interface entre ces
niveaux d'abstraction. Dans ce chapitre nous allons voir son rôle et une
implantation possible. Pour une description plus détaillée, on pourra se référer
à l'ouvrage \cite{tanenbaum}.

Au cours de l'histoire des systèmes informatiques, la manière de les programmer
a beaucoup évolué. Au départ, les programmeurs avaient accès au matériel dans
son intégralité : toute la mémoire pouvait être accédé, toutes les instructions
pouvaient être utilisées.

Néanmoins c'est un peu restrictif : si on est seul à utiliser un système, on est
par définition... seul à pouvoir l'utiliser. Dans la seconde moitié des années
60, sont apparus les premiers systèmes "à temps partagé", permettant à plusieurs
utilisateurs de travailler en même temps.

Permettre l'exécution de plusieurs programmes en même temps est une idée
révolutionnaire, mais elle n'est pas sans difficultés techniques : en effet les
ressources de la machine doivent être aussi partagées entre les utilisateurs et
les programmes. Par exemple, plusieurs programmes vont par exemple utiliser le
CPU les uns à la suite des autres (partage \emph{temporel}) ; et chaque
programme aura à sa disposition une partie de la mémoire principale, ou du
disque dur (partage \emph{spatial}).

Si deux programmes (ou plus) s'exécutent de manière concurrente sur le même
matériel, il faut s'assurer par exemple que les deux s'exécutent à peu près
aussi souvent, ou que l'un ne puisse pas écrire dans la mémoire de l'autre. Ce
sont des rôles du système d'exploitation.

Cela passe donc par un certain bridage des possibilités du programme : plutôt
que de le faire exécuter n'importe quel type d'instruction, il communique avec
le système d'exploitation. Bien que cela ait l'air d'une limitation, c'est aussi
bénéfique pour le programmeur puisque cela permet de définir des abstractions au
niveau du noyau.

Par exemple, si un programmeur veut copier des données depuis un CD-ROM vers la
mémoire principale, il devra interroger le bus SATA, interroger le lecteur sur
la présence d'un disque dans le lecteur, activer le moteur, calculer le numéro
de \emph{frame} des données sur le disque, demander la lecture, puis déclencher
une copie de la mémoire.

Si dans un autre cas il voulait récupérer des données depuis une mémorette USB,
il devrait interroger le bus USB, rechercher le bon numéro de périphérique, le
bon \emph{endpoint} dans celui-ci, lui appliquer une commande de lecture au bon
numéro de bloc, puis copier la mémoire.

Ces deux opérations, bien qu'elles aient le même but (copier de la mémoire
depuis un périphérique amovible), ne sont pas effectuées en pratique de la même
manière. C'est pourquoi le système d'exploitation fournit les notions de
fichier, lecteur, etc : le programmeur n'a plus qu'à utiliser des commandes de
haut niveau (``monter un lecteur'', ``ouvrir un fichier'', ``lire dans un
fichier'') et selon le type de lecteur, le système d'exploitation effectuera les
actions appropriées.

En résumé, un système d'exploitation est l'intermédiaire entre le logiciel et
le matériel, et en particulier assure les rôles suivants :

\begin{itemize}
\item
  Gestion des processus : un système d'exploitation peut permettre
  d'exécuter plusieurs programmes à la fois. Il faut alors orchestrer
  ces différents processus et les séparer en terme de temps et de
  ressources partagées.
\item
  Gestion de la mémoire : chaque processus, en plus du noyau, doit
  disposer d'un espace mémoire différent. C'est-à-dire qu'un processus
  ne doit pas pouvoir interférer avec un autre.
\item
  Gestion des périphériques : le noyau étant le seul code à s'exécuter
  en mode privilégié, c'est lui qui doit communiquer avec les
  périphériques matériels.
\item
  Abstractions : le noyau fournit aux programmes une interface unifiée,
  permettant de stocker des informations de la même manière sur un
  disque dur ou une clef USB (alors que l'accès se déroulera de manière
  très différente en pratique). C'est ici que la notion arbitraire de
  fichier sera définie, par exemple.
\end{itemize}

\section{Architecture Intel}

Dans cette section nous décrivons le fonctionnement d'un processeur utilisant
une architecture Intel 32 bits. Les morceaux de code seront écrit en syntaxe
AT\&T, celle que comprend l'assembleur GNU.

\subsection{Assembleur}

Pour faire des calculs, le processeur est composé de registres, qui sont des
petites zones de mémoire interne, et peut accéder à la mémoire principale.

Les registres principaux sont nommés \eax, \ebx, \ecx, \edx, \esi et \edi. Ils
peuvent être utilisés pour n'importe quel type d'opération, mais certains sont
spécialisés : par exemple il est plus efficace d'utiliser \eax en accumulateur,
ou \ecx en compteur. \esp est le pointeur de pile (\emph{stack pointer}), qui a
un rôle particulier dans les instructions \asminstr{push} et \asminstr{pop}.
\ebp sert de point de repère dans les appels de fonction. \eip contient
l'adresse de l'instruction courante.

Les calculs sont décrits sous forme d'une suite d'instructions. Chaque
instruction est composée d'une mnémonique et d'une liste d'opérandes. Les
mnénoniques (\asminstr{mov}, \asminstr{call}, \asminstr{sub}, etc) définissent
un type d'opération à appliquer sur les opérandes.

Ces opérandes peuvent être de plusieurs types :

\begin{itemize}
\item
  un nombre, noté \$4
\item
  le nom d'un registre, noté \texttt{\%eax}
\item
  une opérande mémoire, c'est à dire le contenu de la mémoire à une
  adresse effective. Cette adresse effective peut être exprimée de
  plusieurs manières :
  \begin{itemize}
  \item
    directement : \texttt{addr}
  \item
    indirectement : \texttt{(\%ecx)}. L'adresse effective est le contenu
    du registre.
  \item
    ``base + déplacement'' : \texttt{4(\%ecx)}. L'adresse effective est
    le contenu du registre plus le déplacement (4 ici).
  \end{itemize}
\end{itemize}

En pratique il y a des modes d'adressage plus complexes, et toutes les
combinaisons ne sont pas possibles, mais ceux-ci suffiront à décrire les
exemples suivants :

\begin{itemize}

\item \texttt{mov src, dst} copie le contenu de \texttt{src} dans \texttt{dst}.

\item \texttt{add src, dst} calcule la somme des contenus de \texttt{src} et
  \texttt{dst} et place ce résultat dans \texttt{dst}.

\item \texttt{push src} place \texttt{src} sur la pile, c'est à dire que cette
  instruction décrémente le pointeur de pile \esp de la taille de \texttt{src},
  puis place \texttt{src} à l'adresse mémoire de la nouvelle valeur \esp.

\item \texttt{pop src} réalise l'opération inverse : elle charge le contenu de
  la mémoire à l'adresse \esp dans \texttt{src} puis incrémente \esp de la
  taille correspondante.

\item \texttt{jmp addr} saute à l'adresse \texttt{addr} : c'est l'équivalent de
  \texttt{mov addr, \%eip}.

\item \texttt{call addr} sert aux appels de fonction : cela revient à
  \texttt{push \%eip} puis \texttt{jmp addr}.

\item \texttt{rest} sert à revenir d'une fonction : c'est l'équivalent de
  \texttt{pop \%eip}.

\end{itemize}

\subsection{Appel de fonction et conventions d'appel}

\begin{figure} % {{{ fig:stackframe
\begin{tikzpicture}
[ stack/.style={draw,shape=rectangle,minimum height=1cm,minimum width=5cm}
, bigbrace/.style={decorate,decoration={brace,amplitude=3mm}}
, >=latex
]

  \path (0,-8.5)   node [stack, minimum height=2cm] (gloc) {Locales de g}
  -- ++(0,1.5) node [stack] (gebp) {Cadre précédent}
  -- ++(0,1)   node [stack] (gret) {Retour de g}
  -- ++(0,1.5) node [stack, minimum height=2cm] (gpar) {Arguments de g}
  -- ++(0,1.5) node [stack] (floc) {Locales de f}
  -- ++(0,1)   node [stack] (febp) {Cadre précédent}
  -- ++(0,1)   node [stack] (fret) {Retour de f}
  -- ++(0,1)   node [stack] (fpar) {Arguments de f};

  \draw [->] (gebp) -- ++(4, 0) |- (febp.355);
  \draw [->] (febp.5) -- ++(1.5, 0) -- ++(0,4);

  \draw [dashed] (fpar.north west) -- ++ (0,2);
  \draw [dashed] (fpar.north east) -- ++ (0,2);
  \draw [dashed] (gloc.south east) -- ++ (0,-2);
  \draw [dashed] (gloc.south west) -- ++ (0,-2);

  \node at (5, 1) (ox) {};

  \node [text width=3cm] at (-0.8, 2)     {Adresses hautes (sens de \asminstr{pop})};
  \node [text width=3cm] at (-0.8, -10.5) {Adresses basses (sens de \asminstr{push})};

  \node [text width=3cm] at (-6, -3) (asmtext) {\tt
    call f\\
    ... \\
    f:  ... \\
    pushl \$2 \\
    pushl \$4 \\
    call g \\
    addl \$8, \%esp \\
    leave \\
    ret \\
    ... \\
    g: push \%ebp \\
    mov \%esp, \%ebp \\
    sub \$8, \%esp \\
    ... \\
    leave \\
    ret
    };

  \draw [->] (gret) -- ++ (-3, 0) |- ($ (asmtext.north west) + (2.5, -3) $);
  \draw [->] (fret) -- ++ (-3, 0) |- ($ (asmtext.north west) + (2.5, -0.5) $);

  \draw [<-] (gebp.west) -- ++ (-0.5,0) node [xshift=-4mm] {\ebp};
  \draw [<-] ($ (gloc.west) + (0,-0.5) $) -- ++ (-0.5,0) node [xshift=-4mm] {\esp};

  \draw [<-] ($ (asmtext.north west) + (2.5, -6.7) $) -- ++ (0.5,0) node [xshift=3mm] {\eip};

  \draw [bigbrace]
        ($ (gpar.north east) + (3, -0.1) $)
     -- node [auto,midway,xshift=3mm] {Cadre de g}
        ($ (gloc.south east) + (3, 0) $);

  \draw [bigbrace]
        ($ (fpar.north east) + (3, 0) $)
     -- node [auto,midway,xshift=3mm] {Cadre de f}
        ($ (floc.south east) + (3, 0.1) $);

\end{tikzpicture}

\caption[Cadres de pile]{ Cadres de pile. }
\label{fig:stackframe}
\end{figure} % }}}

Dans le langage d'assemblage, il n'y a pas de notion de fonction ; mais
\asminstr{call} et \asminstr{ret} permettent de sauvegarder et de restaurer une
adresse de retour, ce qui permet de faire un saut et revenir à l'adresse initiale.
Ce système permet déjà de créer des procédures, c'est-à-dire des fonctions sans
arguments ni valeur de retour.

Pour gérer ceux-ci, il faut que les deux morceaux (appelant et appelé) se
mettent d'accord sur une convention d'appel commune. La convention utilisée sous
GNU/Linux est appelée \emph{cdecl} et possède les caractéristiques suivantes :

\begin{itemize}
\item la valeur de retour d'une fonction est stockée dans \eax
\item \eax, \ecx et \edx peuvent être écrasés sans avoir à les sauvegarder
\item les arguments sont placés sur la pile (et enlevés) par l'appelant. Les
  paramètres sont empilés de droite à gauche.
\end{itemize}

Pour accéder à ses paramètres, une fonction peut donc utiliser un adressage
relatif à \esp. Cela peut fonctionner, mais cela complique les choses si elle
contient aussi des variables locales. En effet, les variables locales sont
placées sur la pile, au dessus des (c'est à dire, empilées après) paramètres,
augmentant le décalage.

Pour simplifier, la pile est organisée en cadres logiques : chaque cadre
correspond à un niveau dans la pile d'appels de fonctions. Si f appelle g, qui
appelle h, il y aura dans l'ordre sur la pile le cadre de f, celui de g puis
celui de h.

Ces cadres sont chainés à l'aide du registre \ebp : à tout moment, \ebp contient
l'adresse du cadre de l'appelant.

Prenons exemple sur la figure~\ref{fig:stackframe} : pour appeler
\texttt{g(4,2)}, \texttt{f} empile les arguments de droite à gauche.
L'instruction \asminstr{call g} empile l'adresse de l'instruction suivante sur
la pile puis saute au début de \texttt{g}.

Au début de la fonction, les trois instructions permettent de sauvegarder
l'ancienne valeur de \ebp, faire pointer \ebp à une endroit fixe dans le cadre
de pile, puis allouer 8 octets de mémoire pour les variables locales.

Dans le corps de la fonction \texttt{g}, on peut donc se référer aux variables
locales par \texttt{-4(\%ebp)}, \texttt{-8(\%ebp)}, etc, et aux arguments par
  \texttt{8(\%ebp)}, \texttt{12(\%ebp)}, etc.

À la fin de la fonction, l'instruction \asminstr{leave} est équivalente à
\texttt{mov \%ebp, \%esp} puis \texttt{pop \%ebp} et permet de défaire le cadre
de pile, laissant l'adresse de retour en haut de pile. Le \asminstr{ret} final
la dépile et y saute.

\subsection{Tâches, niveaux de privilèges}
\label{sec:taches}

Sans mécanisme particulier, le processeur exécuterait uniquement une suite
d'instruction à la fois. Pour lui permettre d'exécuter plusieurs tâches, un
système de partage du temps existe.

À des intervalles de temps réguliers, le système est programmé pour recevoir une
interruption. C'est une condition exceptionnelle (au même titre qu'une division
par zéro) qui fait sauter automatiquement le processeur dans une routine de
traitement d'interruption. Si à cet endroit le code sauvegarde les registres et
restaure un autre ensemble de registres, on peut exécuter plusieurs tâches de
manière entrelacée. Si l'alternance est assez rapide, cela peut donner
l'illusion que les programmes s'exécutent en parallèle. Comme l'interruption
peut survenir à tout moment, on parle de multitâche préemptif.

En plus de cet ordonnancement de processus, l'architecture Intel permet
d'affecter des niveaux de privilège à ces tâches, en restreignant le type
d'instructions exécutables, ou en donnant un accès limité à la mémoire aux
tâches de niveaux moins elevés.

Il y a 4 niveaux de privilèges (nommés aussi \emph{rings}) : le \ring{0} est le
plus privilégié, le \ring{3} le moins privilégié. Dans l'exemple précédent, on
pourrait isoler l'ordonnanceur de processus en le faisant s'exécuter en \ring{0}
alors que les autres tâches seraient en \ring{3}.

\begin{center}\rule{3in}{0.4pt}\end{center}


\begin{figure}
\begin{Verbatim}
+-------------------------------------+
|   Ring 3                            |
|                                     |
|       Programmes utilisateur        |
|                                     |
|                                     |
|   ^ | appels système                |
+---|-v-------------------------------+
|                                     |
|   Ring 0                            |
|                                     |
|       Noyau                         |
|                                     |
|   ^ | instructions privilégiées     |
+---|-v-------------------------------+

   Hardware
\end{Verbatim}

\begin{tikzpicture}

\begin{scope}

  %\path[clip] (0, 0.8) rectangle (5, -0.8);

  \foreach \x/\c in { 4/20, 3/30, 2/40, 1/50, 0/60 }
  \draw[fill=black!\c] ($ (-2,0) + (\x,0) $) circle (2 cm);

\end{scope}

%\draw [<-] (3.5, 0) -- ++(0,-2) node {Applications};
%\draw [<-] (0, 0) -- ++(0,-2) node {Noyau};

\end{tikzpicture}

\caption[Les différents \emph{rings}]{
  Les différents \emph{rings}. Seul le \ring{0} a accès au hardware
  via des instructions privilégiées. Pour accéder aux fonctionnalités du noyau,
  les programmes utilisateur doivent passer par des appels système.}
\label{fig:rings}
\end{figure}

La première protection nécessaire est d'isoler le système d'exploitation
lui-même des programmes qu'il va permettre d'exécuter. Pour ce faire, Le
processeur permet d'exécuter des tâches selon plusieurs niveaux de privilège,
aussi appelés \emph{rings} (figure~\ref{fig:rings}) : du \ring{3}, le
moins privilégié, jusqu'au \ring{0}, le plus privilégié. On peut configurer
le processeur de manière à ce que les instructions privilégiées (accès aux ports
d'entrée/sortie...) ne soient possibles qu'en \ring{0}. Bien que 4 niveaux
soient disponibles, Linux n'utilise que les \emph{rings} 0 et 3 : le noyau
lui-même en \ring{0} et les processus utilisateur en \ring{3}.

\subsection{Mémoire virtuelle}

À partir du moment où plusieurs processus s'exécutent de manière concurrente, un
problème d'isolation se pose : si un processus peut lire dans la mémoire d'un
autre, des informations peuvent fuiter ; et s'il peut y écrire, il peut
détourner l'exécution.

\todo{defs : intégrité, confidentialité}

Le mécanisme de mémoire virtuelle permet de donner à deux tâches une vue
différente de la mémoire : c'est à dire que vue de tâches différentes, une
adresse contiendra une valeur différente.

\begin{figure}
\begin{tikzpicture}
[ bit/.style={ draw
             , shape=rectangle
             , minimum width=4mm
             , minimum height=4mm
             , node distance=4mm
             }
, pstuff/.style={ draw
                , shape=rectangle
                , node distance=1cm
                , text width=2cm
                , text centered
                }
, >=latex
]

\node [bit] (b31) {};

% generate this ?
\node [bit, right of=b31] (b30) {}; \node [bit, right of=b30] (b29) {}; \node [bit, right of=b29] (b28) {};
\node [bit, right of=b28] (b27) {}; \node [bit, right of=b27] (b26) {}; \node [bit, right of=b26] (b25) {};
\node [bit, right of=b25] (b24) {}; \node [bit, right of=b24] (b23) {}; \node [bit, right of=b23] (b22) {};
\node [bit, right of=b22] (b21) {}; \node [bit, right of=b21] (b20) {}; \node [bit, right of=b20] (b19) {};
\node [bit, right of=b19] (b18) {}; \node [bit, right of=b18] (b17) {}; \node [bit, right of=b17] (b16) {};
\node [bit, right of=b16] (b15) {}; \node [bit, right of=b15] (b14) {}; \node [bit, right of=b14] (b13) {};
\node [bit, right of=b13] (b12) {}; \node [bit, right of=b12] (b11) {}; \node [bit, right of=b11] (b10) {};
\node [bit, right of=b10] (b9) {}; \node [bit, right of=b9] (b8) {}; \node [bit, right of=b8] (b7) {};
\node [bit, right of=b7] (b6) {}; \node [bit, right of=b6] (b5) {}; \node [bit, right of=b5] (b4) {};
\node [bit, right of=b4] (b3) {}; \node [bit, right of=b3] (b2) {}; \node [bit, right of=b2] (b1) {};
\node [bit, right of=b1] (b0) {};

\foreach \b in {0,11,12,21,22,31}
    \node [above of=b\b, node distance=5mm] {\small \b};

\newcommand{\mbrace}[3]{
    \draw [bigbrace]                      ($ (b#1.south east) + (-0.1,-0.3) $)
    -- node [midway, yshift=-2mm] (#3) {} ($ (b#2.south west) +  (0.1,-0.3) $);
} % mbrace

\mbrace{0}{11}{bpp};
\mbrace{12}{21}{bpt};
\mbrace{22}{31}{bpd};

\node[pstuff, below of=bpd] (pd) {Répertoire de pages};
\node[pstuff, below of=bpt] (pt) {Table de pages};
\node[pstuff, below of=bpp] (pp) {Page de 4 Kio};

\node[left of=pd, node distance=2cm] (cr) {\crtrois};

\draw [->] (bpp) -- (pp);
\draw [->] (bpt) -- (pt);
\draw [->] (bpd) -- (pd);

\draw [->] (cr) -- (pd);
\draw [->] (pd) -- (pt);
\draw [->] (pt) -- (pp);
\draw [->] (pp) -- ++ (3,0);

\end{tikzpicture}

\caption{Implémentation de la mémoire virtuelle}
\label{fig:pagetables}
\end{figure}

Ce mécanisme est contrôlé par valeur du registre \crtrois : les 10 premiers bits
d'une adresse sont un index dans le répertoire de pages qui commence à l'adresse contenue dans
\crtrois. À cet index, se trouve l'adresse d'une table de pages. Les 10 bits
suivants de l'adresse sont un index dans cette page, donnant l'adresse d'une
page de 4 kibi-octets (figure~\ref{fig:pagetables}).

\begin{figure} % fig:memoire-virtuelle {{{
\centering
\begin{tikzpicture}
  [every node/.style={node distance=2cm}]

  \node (u1) {};

  \node[right of=u1] (phy) {};

  \node[right of=phy] (u2) {};

  \draw (u1) rectangle ++(1,-6);
  \draw (phy) rectangle ++(1,-6);
  \draw (u2) rectangle ++(1,-6);
\end{tikzpicture}
\caption{Mécanisme de mémoire virtuelle.}
\label{fig:memoire-virtuelle}
\end{figure} % }}}

En ce qui concerne la mémoire, les différentes tâches ont une vision différente
de la mémoire physique : c'est-à-dire que deux tâches lisant à une même adresse
peuvent avoir un résultat différent. C'est le concept de mémoire virtuelle
(fig~\ref{fig:memoire-virtuelle}).

\begin{figure} % fig:memmap {{{
\centering
\fbox{
  \begin{tikzpicture}
  [scale=0.7
  ,user/.style={fill=black!20}
  ,kernel/.style={fill=black!70}
  ]

  % Memory zone
  %
  % #1 - start
  % #2 - end
  % #3 - color
  \newcommand{\mzone}[3]{
    \path[#3] (#1,0) rectangle (#2,1);
  }

  % Address label
  %
  % #1 - x position
  % #2 - text
  \newcommand{\alabel}[2]{
    \path (#1,1) -- ++(0,0.3) node [pos=1] {\small \tt #2};

  }

  % exec
  \mzone{0.5}{1}{user}

  % lib
  \mzone{2.5}{3.2}{user}

  % stack
  \mzone{3.7}{4}{user}

  % stack
  \mzone{5}{5.5}{user}

  % kernel
  \mzone{6}{8}{kernel}

  % contour
  \draw (0,0) rectangle (8,1);

  \alabel{0}{0}
  \alabel{6}{3 Go}
  \alabel{8}{4 Go}

\end{tikzpicture}

}

\caption[Espace d'adressage d'un processus]{L'espace d'adressage d'un processus.
En gris clair, les zones accessibles à tous les niveaux de privilèges : code du
programme, bibliothèques, tas, pile. En gris foncé, la mémoire du noyau,
réservée au mode privilégié.}

\label{fig:memmap}
\end{figure}
% }}}

\section{Cas de Linux}

Dans cette section, nous allons voir comment ses mécanismes sont implantés dans
le noyau Linux. Une description plus détaillée pourra être trouvée dans
\cite{UnderstandingTheLinuxKernel} par exemple.

Deux rings sont utilisés : en \ring{0}, le code noyau et en \ring{3}, le code
utilisateur.

Une notion de tâche similaire à celle décrite dans la section~\ref{sec:taches}
existe : elles s'exécutent l'une après l'autre, le changement s'effectuant sur
interruptions.

Pour faire appel aux services du noyau, le code utilisateur peut faire appel à
des appels systèmes, qui sont des fonctions exécutées par le noyau. Chaque tâche
doit donc avoir deux piles : une pile "unitisateur" qui sert pour l'application
elle-même, et une pile "noyau" qui sert aux appels système.

\subsection{Appels système}

Les programmes utilisateur s'exécutant en \ring{3}, ils ne peuvent pas
contenir d'instructions privilégiées, et donc ne peuvent pas accéder directement
au matériel (c'était le but !). Pour qu'ils puissent interagir avec le système
(afficher une sortie, écrire sur le disque...), le mécanisme des appels système
est nécessaire. Il s'agit d'une interface de haut niveau entre les \emph{rings}
3 et 0. Du point de vue du programmeur, il s'agit d'un ensemble de fonctions C
``magiques'' qui font appel au système d'exploitation pour effectuer des
opérations.

Prenons le cas de l'appel \texttt{getpid}, qui retourne le numéro de processus
courant. La bibliothèque C fournit une fonction du même nom :

\begin{Verbatim}
pid_t getpid(void);
\end{Verbatim}

A priori, rien de différent d'une fonction implantée directement en C. Par un
processus détaillé ci-après, cette fonction va invoquer la fonction, suivante,
définie dans le noyau (\texttt{kernel/timer.c}) :

\begin{Verbatim}
SYSCALL_DEFINE0(getpid)
{
        return task_tgid_vnr(current);
}
\end{Verbatim}



Le mécanisme de couplage entre ces deux fonctions est le suivant. Une
explication plus détaillée est disponible dans la documentation d'Intel
\cite{intelsys}.

Il y a bien une fonction \texttt{getpid} présente dans la bibliothèque C du
système. C'est la fonction qui est directement appelée par le programme. Cette
fonction commence par placer le numéro de l'appel système (noté
\texttt{\_\_NR\_getpid}, valant 20 ici) dans \eax, puis les arguments éventuels
dans les registres (\ebx, \ecx, \edx, \esi puis \edi). Une interruption
logicielle est ensuite déclenchée (\texttt{int 0x80}) : étant donné la
configuration du processeur, elle sera traitée en \ring{0}, à un point
d'entrée prédéfini (\texttt{arch/x86/kernel/entry\_32.S},
\texttt{ENTRY(system\_call)}).

L'exécution reprend donc en \ring{0}, avec une pile propre au processus.
Les valeurs des registres ont été préservées : elles sont mises sur la pile du
noyau. Le noyau examine le numéro d'appel système (toujours dans \eax) et
appelle la fonction correspondante dans le tableau \texttt{sys\_call\_table} :
les arguments sont en place sur la pile, comme dans le cas d'un appel de
fonction ``classique''. La fonction, nommée \texttt{sys\_getpid}, peut donc être
écrite en C.

Au retour de la fonction, la valeur de retour est placée à la place de \eax là
où les registres ont été sauvegardés sur la pile noyau. L'instruction
\texttt{iret} permet de restaurer les registres et de repasser en mode
utilisateur, juste après l'interruption. La fonction de la bibliothèque C peut
alors retourner au programme appelant.

% vim: spelllang=fr


%\chapter{Typage}

%Dans ce chapitre, nous explorons la notion de type dans les langages de
programmation. Tout d'abord pourquoi elle existe et en quoi elle aide à rendre
les programmes plus sûrs. Il y a autant de système de types que de langages de
programmation, donc nous présenterons ensuite une taxonomie de ces systèmes en
les regroupant par caractéristiques communes. Cette classification sera appuyée
par des exemples de code C\cite{AnsiC,KandR}, OCaml\cite{ocamlManual,DAOC},
Haskell\cite{haskell98,rwh}, Python\cite{pythonSite} et
Perl\cite{perlCamelBook}.

\section{Présentation et but}

Nous avons vu dans le chapitre~\ref{cha:os} qu'au niveau du langage machine, les
seules données qu'un ordinateur manipule sont des nombres. Selon les opérations
effectuées, ils seront interprétés comme des entiers, des adresses mémoires, ou
des caractères. Pourtant il est clair que certaines opérations n'ont pas de
sens: par exemple, ajouter deux adresses, ou déréférencer le résultat d'une
division sont des comportements qu'on voudrait pouvoir empêcher.

En un mot, le but du typage est de classifier les objets et de restreindre les
opérations possibles selon la classe d'un objet : "ne pas ajouter des pommes et
des oranges". Le modèle qui permet cette classification est appelé \emph{système
de types} et est en général constitué d'un ensemble de \emph{règles de typage},
comme "un entier plus un entier égale un entier".

\section{Taxonomie}

La définition d'un langage de programmation introduit la plupart du temps celle
d'un système de types. Il y a donc de nombreux systèmes de types différents,
dont nous pouvons donner une classification sommaire.

\subsection{Dynamique, statique, mixte}

Il y a deux grandes familles de systèmes de types, selon quand se fait la
vérification de types. On peut en effet l'effectuer au moment de l'exécution, ou
au contraire prévenir les erreurs à l'exécution en la faisant au moment de la
compilation (ou avant l'interprétation).

\subsubsection{Typage dynamique}

La première est le typage \emph{dynamique}. Pour différencier les différents
types de données, on ajoute une étiquette à chaque valeur. Dans tout le
programme, on ne manipulera que des valeurs étiquetées, c'est à dire des couples
(donnée, nom de type). Si on veut réaliser l'opération $(\hexa{00000001}, Int) +
(\hexa{0000f000}, Int)$, on vérifie tout d'abord qu'on peut réaliser l'opération
$+$ entre deux $Int$. Ensuite on réalise l'opération elle même, qu'on étiquette
avec le type du résultat : $(\hexa{0000f001}, Int)$. Si au contraire on tente
d'ajouter deux adresses $(\hexa{2e8d5a90}, Addr) + (\hexa{76a5e0ec}, Addr)$, la
vérification échoue et l'opération s'arrête avec une erreur.

\begin{figure}
  \insertcode{typage-dynamique.py}
  \caption{Session Python présentant le typage dynamique}
  \label{fig:typage-dynamique}
\end{figure}

La figure~\ref{fig:typage-dynamique} est une session interactive Python qui
illustre le typage dynamique. Chaque variable, en plus de sa valeur, possède une
étiquette qui permet d'identifier le type de celle-ci. On peut accéder au type
d'une valeur \texttt{x} avec la construction \texttt{type(x)}.

Au moment de réaliser une opération comme \texttt{+}, l'interpréteur Python
vérifie les types des opérandes : s'ils sont compatibles, il créé une valeur de
résultat, et sinon il lève une exception.

Comme l'implémentation elle-même des fonction a accès aux informations de type,
elle peut faire des traitements particuliers. Par exemple, pour l'addition de
\texttt{a} (de type entier) et de \texttt{c} (de type flottant), la fonction
d'addition va d'abord convertir \texttt{a} en flottant, puis réaliser l'addition
dans le domaine des flottants.

\subsubsection{Typage statique}

Le typage dynamique est simple à comprendre puisque toute les vérifications se
font dans la sémantique dynamique (ie, à l'exécution). C'est à double tranchant:
d'une part, cela permet d'être plus flexible, mais d'autre part, cela permet à
des programmes incorrects d'être exécutés.

On peut lire le code source d'un programme et essayer de "deviner" quels seront
les types des différentes expressions. Dans certains cas, cela n'est pas
possible (fig~\ref{fig:nontypable}) ; mais lorsqu'on peut conclure cela élimine
la nécessité de faire les tests de type dynamiques car on a réalisé le typage
\emph{statiquement}.

Bien sûr, deviner n'est pas suffisant : il faut formaliser cette analyse. Dans
le cas dynamique, ce sont les fonctions elles-mêmes qui réalisent les tests de
type et qui appliquent des règles comme "si les arguments ont pour type
\texttt{int} alors la valeur de retour a pour type \texttt{int}" : la fonction
qui réalise ce test sur les valeurs. Dans le cas statique, c'est le compilateur
(ou l'interpréteur) qui réalise ce test sur les expressions non évaluées. En
appliquant de proche en proche un ensemble de règles (liées uniquement au
langage de programmation), on finit par associer à chaque sous-expression du
programme un type.

Benjamin Pierce résume cette approche dans cette définition très générale :

\begin{definition}[Système de types]
Un système de types est une méthode syntaxique tractable qui vise à prouver
l'absence de certains comportements de programmes en classifiant leurs phrases
selon le genre de valeurs qu'elles produisent. \cite{TAPL}
\end{definition}

\begin{figure}
  \insertcode{non-typable.py}
  \caption{Fonction Python non typable statiquement.}
  \label{fig:nontypable}
\end{figure}

A première vue, cela semble moins puissant que le typage dynamique : en effet,
il existe des programmes qui s'exécuteront sans erreur de type mais sur lesquels
le typage statique ne peut s'appliquer. Dans la figure~\ref{fig:nontypable}, on
peut voir par une simple analyse de cas que si on fournit un booléen à
\texttt{f}, elle retourne un entier. Mais selon la valeur de \texttt{b}, la
variable \texttt{x} contiendra une valeur de type entier ou fonction.

Même si cet exemple est construit artificiellement, il illustre le problème
suivant : les types statiques demandent un certain effort et au programmeur et
au compilateur. Mais Dans le cas où le typage statique est possible, les
garanties sont importantes : les valeurs portées par une variable auront
toujours le même type. Par voie de conséquence, la vérification dynamique de
types réussira toujours, et on peut la supprimer. Il est également possible
de supprimer toutes les étiquettes de typage : on parle de \emph{type erasure}.
Une conséquence heureuse de cette suppression est que l'exécution de ce
programme se fera de manière plus rapide.

Connaître les types à la compilation permet aussi de réaliser plus
d'optimisations. Par exemple, en Python, considérons l'expression \texttt{y = x
- x}. Sans information sur le type de \texttt{x}, aucune simplification n'est
possible : l'implémentation de la différence sur ce type est une fonction
quelconque, sans propriétés particulières \emph{a priori}. Si au contraire, on
sait que \texttt{x} est un entier, on peut en déduire que \texttt{y = 0}, sans
réaliser la soustraction (si c'était la seule utilisation de \texttt{x}, le
calcul de \texttt{x} aurait alors pu être éliminé).

\subsection{Fort, faible, sound}

\begin{figure}
  \insertcode{cast.java}
  \caption{Transtypage en Java}
  \label{fig:javacast}
\end{figure}

Si un système de types statique permet d'éliminer totalement la nécessité de
réaliser des tests de typage, on dit qu'il est \emph{fort}. Mais ce n'est que
rarement le cas. En effet, il peut y avoir des constructions au sein du langage
qui permettent de contourner le système de types, comme un opérateur de
transtypage\ref{fig:javacast}. À l'exécution, une erreur de types est levée :

\begin{Verbatim}
Exception in thread "main" java.lang.ClassCastException:
    java.lang.Integer cannot be cast to java.lang.Float
        at Cast.main(Cast.java:5)
\end{Verbatim}

\clearpage

\wip{}

\subsection{Polymorphisme}

Dans le cas du typage statique, restreindre une opération à un seul type de
données peut être assez restrictif.

Par exemple, quel doit être le type d'une fonction qui trie un tableau ?

\begin{figure}
\centering
\begin{tikzpicture}
  [ ptype/.style={draw=black,shape=rectangle}
  , child/.style={ptype, xshift=2cm}
  , gen0/.style={       node distance=2cm}
  , gen1/.style={child, node distance=1cm}
  , gen2/.style={child, node distance=5mm}
  ]

  \node[ptype]                     (mono) {Mono};
  \node[ptype,below of=mono, gen0] (poly) {Poly};
  \node[above right of=poly, gen1] (univ)  {Universel};
  \node[above right of=univ, gen2] (param) {Paramétrique};
  \node[below right of=univ, gen2] (inclu) {Inclusion};
  \node[below right of=poly, gen1] (adho)  {Ad-hoc};
  \node[above right of=adho, gen2] (over)  {Surcharge};
  \node[below right of=adho, gen2] (coer)  {Coercition};

  \draw (poly) -- (univ);
  \draw (poly) -- (adho);

  \draw (univ) -- (param);
  \draw (univ) -- (inclu);

  \draw (adho) -- (over);
  \draw (adho) -- (coer);

      %node {Mono}
      %node {Poly}
          %node {Universal}
              %node {Param}
              %node {Inclusion}
          %node {Ad-Hoc}
              %node {Surcharge}
              %node {Coercition}

\end{tikzpicture}

\caption{Les différents types de polymorphisme.}
\label{fig:types-de-polymorphisme}
\end{figure}

\subsubsection{Monomorphisme}

Une première solution peut être de forcer des types concrets, c'est à dire
qu'une même fonction ne pourra s'appliquer qu'à un seul type de données.

Il est confortable pour le programmeur de n'avoir à écrire un algorithme qu'une
seule fois, indépendamment du type des éléments considérés.

Il existe deux grandes classes de systèmes de types introduisant du
polymorphisme.

\subsubsection{Polymorphisme universel}

% = {paramétrique, par inclusion}

Le polymorphisme est dit universel si toute fonction générique peut s'appliquer
à n'importe quel type.

\subsubsection{Polymorphisme ad-hoc}

% = {par surcharge, par coercition}

Le polymorphisme est \emph{ad-hoc} si les fonctions génériques ne peuvent
s'appliquer qu'à un ensemble de types prédéfini.

\subsubsection{Polymorphisme paramétrique}

\todo{Historique + citer le papier de Milner sur le polymorphisme}
\cite{Milner78}

\cite{PascalNoEscape}

\begin{figure}
  \insertcode{listappend.ml}
  \caption{Fonction de concaténation de listes en OCaml.}
  \label{fig:listappend}
\end{figure}

La fonction de la figure~\ref{fig:listappend} n'opère que sur la structure du type
liste (en utilisant ses constructeurs \texttt{{[}{]}} et \listcons ainsi que
le filtrage) : les éléments de \texttt{lx} et \texttt{ly} ne sont pas manipulés
à part pour les transférer dans le résultat.

Moralement, cette fonction est donc indépendante du type de données contenu dans
la liste : elle pourra agir sur des listes de n'importe quel type d'élément.

Plutôt qu'un type, on peut lui donner le \emph{schéma de types} suivant :

\[
  \textrm{append} : \forall a . a \textrm{list}
                             -> a \textrm{list}
                             -> a \textrm{list}
\]

C'est à dire que \texttt{append} peut être utilisé avec n'importe quel type
concret \texttt{a} en substituant les variables quantifiées (on parle d'
\emph{instanciation}).

\subsubsection{Polymorphisme par sous-typage}

\todo{héritage,sous-typage,classe,méthode,héritage multiple,late binding,Liskov}

Certains langages définissent la notion de sous-typage. C'est une relation
d'ordre partiel sur les types, qui modélise la relation "est un". Chaque
sous-classe peut redéfinir le comportement de chaque méthode de ses
superclasses.

\subsubsection{Polymorphisme par surcharge}

Considérons l'opération d'addition : \texttt{+}. On peut considérer que certains
types l'implémentent, et pas d'autres : ajouter deux flottants ou deux entiers a
du sens, mais pas ajouter deux pointeurs.

On dira que \texttt{+} est \emph{surchargé}. À chaque site d'appel, il faudra
\emph{résoudre la surcharge} pour déterminer quelle fonction appeler.

Cela rend l'inférence de types \todo{introduire l'inférence plus haut}
impossible dans le cas général, puisque certaines constructions sont ambigües.

\begin{figure}
  \insertcode{showread.hs}
  \caption{Cas d'ambigüité avec de la surcharge ad-hoc.}
  \label{fig:showread}
\end{figure}

Dans le code Haskell de la figure~\ref{fig:showread}, \texttt{show} peut
s'appliquer à toutes les valeurs de types "affichables" et renvoie une
représentation textuelle. \texttt{read} réalise le contraire avec les types
"lisibles".

Lorsqu'on compose ces deux fonctions, le type de la valeur intermédiaire est
capital puisqu'il détermine les instances de \texttt{show} et \texttt{read} à
utiliser.

\subsubsection{Polymorphisme par coercition}

\subsubsection{Polymorphisme d'ordre supérieur}

\begin{verbatim}
g f = (f true, f 2)
\end{verbatim}

\[
g : (\forall a . a -> a) -> (bool * int)
\]

Pas inférable (annotations nécessaires).

\subsection{Expressivité, garanties, types dépendants}

\section{Exemples}

\subsection{Faible dynamique : Perl}
\subsection{Faible statique : C}
\subsection{Fort dynamique : Python}
\subsection{Fort statique : OCaml}
\subsection{Fort statique à effets typés : Haskell}
\subsection{Theorem prover : Coq}

% vim: spelllang=fr


\chapter{État de l'art}

\begin{itemize}
\item
  Analyse légère :
\item
  Sparse \cite{TorvaldsSparse}
\item
  Coccinelle \cite{coccinelle09} \cite{coccinelle11} \todo{lire coccinelle09}
\item
  Analyse de valeurs :
\item
  IA :
  \begin{itemize}
  \item
    Cousot 77 \cite{Cousot77}
  \item
    Cousot 92 \cite{Cousot92-1}
  \item
    CGS \cite{cgs}
  \end{itemize}
\item
  Typage fort
\item
  \cite{TAPL}
\item
  Qualifiers in Haskell :
  \begin{itemize}
  \item
    \cite{lightweight-static-capabilities}
  \item
    \cite{LZ06a}
  \end{itemize}
\item
  Qualificateurs de type :
\item
  CQual \cite{pldi99,usenix01,pldi02,cquk-usenix04,toplas-quals}
\end{itemize}


\part{Typage statique de langages impératifs}

\chapter{Sémantique d'un langage impératif}

\label{cha:lang}
On définit ici un langage impératif.

\section{Syntaxe}

Donnée dans la figure~\ref{fig:stx}.

\begin{figure}

  \gramlr{Constantes}{
    \begin{align*}
      c \gramisa & i & \textrm{Entier}
      \\ \gramor & f & \textrm{Flottant}
    \end{align*}
  }

  \gramlr{Expressions}{
    \begin{align*}
      e \gramisa & c               & \textrm{Constante}
      \\ \gramor & x               & \textrm{Variable}
      \\ \gramor & \opsymb~e       & \textrm{Opération unaire}
      \\ \gramor & e~\opsymb~e     & \textrm{Opération binaire}
      \\ \gramor & \& x            & \textrm{Pointeur}
      \\ \gramor & * e             & \textrm{Déréférencement}
      \\ \gramor & f (e_1, …, e_n) & \textrm{Appel de fonction}
    \end{align*}
  }

  \gramlr{Instructions}{
    \begin{align*}
      i \gramisa & \iPass        & \textrm{Instruction vide}
      \\ \gramor & i;i           & \textrm{Séquence}
      \\ \gramor & x ← e         & \textrm{Affectation}
      \\ \gramor & \iIf{e}{i}{i} & \textrm{Alternative}
      \\ \gramor & \iWhile{e}{i} & \textrm{Boucle}
      \\ \gramor & \iReturn{e}   & \textrm{Retour de fonction}
    \end{align*}
  }

  \gramlr{Fonctions}{
    \begin{align*}
      f \gramisa & (i, (x_1, …, x_n)) & \textrm{Corps et arguments}
    \end{align*}
  }

  \caption{Syntaxe}
  \label{fig:stx}
\end{figure}

\section{Sémantique}

On définit dans la suite une sémantique pour le langage défini dans la
figure~\ref{fig:stx}.

\section{État mémoire}

Un programme, pour s'exécuter, dispose de mémoire organisée en variables. Plus
précisément, lors de l'exécution de chaque fonction, un cadre de pile est créé,
contenant un ensemble de variables accessibles par cette fonction. Il y en outre
un ensemble de variables globales accessibles depuis n'importe quel point du
programme.

\begin{figure}

  \gramlr{Valeurs}{
    \begin{align*}
      v \gramisa & c & \textrm{Constante}
      \\ \gramor & a & \textrm{Adresse}
    \end{align*}
  }

  \gramlr{Adresses}{
    \begin{align*}
      a \gramisa & (n, x) & \textrm{Variable locale}
      \\ \gramor & x      & \textrm{Variable globale}
    \end{align*}
  }

  \gramlr{Pile}{
    \begin{align*}
      s \gramisa & []                                   & \textrm{Pile vide}
      \\ \gramor & \{ x_1 ↦ v_1 ; … ; x_n ↦ v_n \} :: s & \textrm{Ajout d'un cadre}
    \end{align*}
  }

  \gramlr{État mémoire}{
    \begin{align*}
      m \gramisa & (s, \{ x_1 ↦ v_1 ; … ; x_n ↦ v_n \}) & \textrm{Pile, globales}
    \end{align*}
  }

  \caption{Interpréteur}
  \label{fig:interp}
\end{figure}

Chaque fonction peut donc accéder aux variables du dernier cadre de pile, ainsi
qu'aux variables globales.

\[
  (x, v) ∈ m
\]

\[
  m[x ↦ v]
\]

\begin{definition}{Manipulations de pile}

  Empilement / dépilement

  \[
    \mathrm{Push}(s, (x_1, v_1), …, (x_n, v_n)) = \{ x_1 ↦ v_1 ; … ; x_n ↦ v_n \}~::~s
  \]

  \[
    \mathrm{Pop}(f~::~s) = s
  \]

  + extension aux $m$.

\end{definition}

\begin{definition}{Recherche de variable}

  \[
    \textrm{VarLookup} (x, (s, g)) =
      \begin{cases}
        x      & \mbox{si } x ∈ g \\
      (0, x)   & \mbox{si } x ∉ g \mbox{ et } s = f::[]
                 \mbox{ avec } x ∈ f \\
      (n+1, x) & \mbox{si } x ∉ g \mbox{ et } s = f::s'
                 \mbox{ avec } \textrm{VarLookup} (x, (s, g)) = (n, x)
      \end{cases}
  \]

\end{definition}

\section{Expressions}

Réduction (grand pas) :

\[
  (m, e) → (m', v)
\]

Constante :

$\widehat{c}$ désigne la valeur (sémantique) de la constante (syntaxique) $c$.

\begin{mathpar}
  \irule{Exp-Cst}
    { }
    {(m, c) → (m, \widehat{c})}
\end{mathpar}

Cas variable :

\begin{mathpar}
  \irule{Exp-Var}
  {(x, v) ∈ m}
  {(m, x) → (m, v)}
\end{mathpar}

Cas op :

\begin{mathpar}
  \irule{Exp-UnOp}
  {(m, e) → (m', v)}
  {(m, \opsymb~e) → (m', \widehat{\opsymb}~v)}
\end{mathpar}

À chaque opérateur $\opsymb$, correspond un opérateur sémantique
$\widehat{\opsymb}$ qui agit sur les valeurs. Par exemple, l'opérateur
$\widehat{+}$ est l'addition classique entre entiers.

\begin{mathpar}
  \irule{Exp-BinOp}
  {(m, e_1) → (m', v_1) \\
   (m', e_2) → (m'', v_2)
  }
  {(m, e_1~\opsymb~e_2) → (m'', v_1~\widehat{\opsymb}~v_2)}
\end{mathpar}

Adresse :

\begin{mathpar}
  \irule{Exp-AddrOf}
    {a = \textrm{VarLookup} (x, m)}
    {(m, \&x) → (m, a)}
\end{mathpar}

Déréférencement :

\begin{mathpar}
  \irule{Exp-AddrOf}
    { (m, e) → (m', a) \\
      (a, v) ∈ m'
    }
    {
      (m, *e) → (m', v)
    }
\end{mathpar}

Appel de fonction :

\[
  \irule{Call}
  {(m  , e_1) → (m_1, v_1) \\
   (m_1, e_2) → (m_2, v_2) \\
   … \\
   (m_{n-1}, e_n) → (m_n, v_n) \\
   f = (i, arg_1, …, arg_n) \\
   m' = \mathrm{Push}(m_n, (arg_1, v_1), …, (arg_n, v_n)) \\
   (m', i) → (m'', \iReturn{v}) \\
   m''' = \mathrm{Pop}(m'')
  }
  {(m, f(e_1, …, e_n)) → (m''', v)}
\]

\section{Instructions}

(attention au pas)

\[
  \irule{Set}
  {
    (m, e) → (m', v) \\
    m'' = m' [ x ↦ v ]
  }
  {(m, x ← e) → (m'', pass)}
\]

\begin{itemize}
\item
  ajout des adresses
\item
  declaration
\item
  champs de records
\end{itemize}


\chapter{Typage}

\label{cha:typbase}
\section{Expressions}

\begin{mathpar}

  \irule{Cst-Int}
    { }
    { Γ ⊢ i : \tInt}

  \irule{Cst-Float}
    { }
    { Γ ⊢ i : \tInt}

  % TODO
  \irule{Cst-Null}
    { }
    { Γ ⊢ \eNull : ?}

  \irule{Lv-Var}
    { (x, t) ∈ Γ }
    { Γ ⊢ x : t }

  \irule{Lv-Deref}
    { Γ ⊢ x : t* }
    { Γ ⊢ *x : t }

  % TODO
  \irule{Lv-Field}
    { }
    { Γ ⊢ x.f : t }

  \irule{Lv-Index}
    { Γ ⊢ e : \tInt \\
      Γ ⊢ x : t[]
    }
    { Γ ⊢ x[e] : t }

  % TODO opérateurs

  \irule{Addr}
    { Γ ⊢ lv : t }
    { Γ ⊢ \&lv : t* }

  \irule{Set}
    { Γ ⊢ lv : t \\
      Γ ⊢ e : t
    }
    { Γ ⊢ lv ← e : t }

  \irule{Struct}
    { ∀ i, 1 ≤ i ≤ n ⇒ Γ ⊢ e_i : t_i }
    { Γ ⊢ \{ f_1 : e_1 ;
         … ; f_n : e_n \}
        : \{ f_1 : t_1 ;
         … ; f_n : t_n \}
    }

  \irule{Array}
    { ∀ i, 1 ≤ i ≤ n ⇒ Γ ⊢ e_i : t }
    { Γ ⊢ \{ e_1 ;
         … ; e_n \}
        : t[]
    }

  \irule{Call}
    { ∀ i, 1 ≤ i ≤ n ⇒ Γ ⊢ e_i : t_i \\
      Γ ⊢ f : ( t_1 ,
            … , t_n ) → t
    }
    { Γ ⊢ f( e_1 ,
         … , e_n )
        : t
    }

\end{mathpar}

\section{Instructions}

\begin{mathpar}

  \irule{Pass}
    { }
    {Γ ⊢ \iPass}

  \irule{Seq}{
    Γ ⊢ i_1 \\
    Γ ⊢ i_2
  }{
    Γ ⊢ i_1;i_2
  }

  \irule{Exp}
    { Γ ⊢ e : t }
    { Γ ⊢ e }

  \irule{If}
    { Γ ⊢ e : \tInt \\
      Γ ⊢ i_1 \\
      Γ ⊢ i_2
    }
    { Γ ⊢ \iIf{e}{i_1}{i_2} }

  \irule{While}
    { Γ ⊢ e : \tInt \\
      Γ ⊢ i
    }
    { Γ ⊢ \iWhile{e}{i} }

  % TODO comment le détecter pour typer la fonction ? faire une fausse lvalue ?
  \irule{Return}
    { Γ ⊢ e : t }
    { Γ ⊢ \iReturn{e} }

\end{mathpar}

\section{Fonctions}

\begin{mathpar}

  \irule{Fun}
    { ?? }
    { Γ ⊢ f : (t_1, …, t_n) → t }

\end{mathpar}

\section{Programme}

\begin{mathpar}
  \irule{Prog}
    { ∀ i, 1 ≤ i ≤ p ⇒ Γ e_i : t_i
    }
    {
      Γ ⊢ ((f_1, …, f_n), ((x_1, e_1), …, (x_p, e_p)))
    }
\end{mathpar}

\begin{center}\rule{3in}{0.4pt}\end{center}

\section*{TODO}

\begin{itemize}
\item
  typer les opérateurs
\item
  utiliser des phrases toplevel
\item
  résoudre le problème du return
\end{itemize}


%\chapter{Sémantique statique}

%\label{cha:typbase}
%Ici nous enrichissons le langage défini dans le chapitre~\ref{cha:lang} d'un
système de types. Celui-ci permet d'obtenir plus de garanties que celui de C tel
que décrit dans \cite{AnsiC}.

Il permet le polymorphisme sur les types pointeurs, permettant par exemple de
typer :

\begin{mathpar}
⊢ \text{memcpy} : ∀ a . (a^*, a^*, \tInt) \rightarrow \tVoid
\end{mathpar}

\section{Règles de typage}

\subsection{Types}

Dans cette section, on définit la notion de programme bien typé. L'analyse par
typage permet de vérifier qu'à chaque expression on peut associer un type, et ce
de manière cohérente entre plusieurs utilisations d'une variable.

Les types des valeurs sont :

\gramlr{Types}{
\begin{align*}
τ  \gramisa & \tInt, \tFloat, \tVoid        & \textrm{Constante}
\\ \gramor  &  a                            & \textrm{Variable}
\\ \gramor  & (τ_1, …, τ_n) \rightarrow τ_r & \textrm{Fonction}
\\ \gramor  & [τ]                           & \textrm{Tableau}
\\ \gramor  & τ*                            & \textrm{Pointeur}
\\ \gramor  & \{ f_1:τ_1
              ,    …
              , f_n:τ_n \}                  & \textrm{Structure}
\end{align*}
}

L'ensemble des types possibles (défini inductivement ci-dessus) sera noté
\textsc{Typ}, et l'ensemble des variables de type par \textsc{VarTyp}.

\subsection{Schémas de type}

On va associer à chaque variable globale un type. Mais faire de même pourrait
être trop restrictif. En effet, une fonction comme memcpy peut être utilisée
pour copier des tableaux d'entiers, mais aussi de flottants. On va donc associer
un schéma de types à chaque fonction.

\gramlr{Schémas}{
\begin{align*}
σ \gramisa & ∀ \vec{a} . τ
\end{align*}
}

Un schéma de types correspond à un ensemble de types. Prenons l'exemple de la
fonction identité : elle a pour schéma de types $∀ a . a \rightarrow a$, ce qui
signifie que pour chaque type $τ$, on peut l'utiliser avec le type $τ
\rightarrow τ$. Plus précisément, cela veut dire que puisque $a$ est quantifiée,
on peut le substituer par n'importe quel autre type.

\begin{definition}[Substitution]
Une substitution est une fonction partielle de \textsc{VarTyp} dans
\textsc{Typ}. Elle sera notée par exemple $s = \{ a ↦ \tInt, b ↦ (\tFloat
\rightarrow \tInt)\}$.

Le domaine de définition d'une substition est noté $\textrm{Dom}(s)$.

On définit aussi l'application d'une substitution sur un type quelconque : si
$s$ est une substitution, $\overline{s}$ est son extension définie par :

\begin{align*}
\overline{s}(a)   =~& s(a) & \textrm{si $a$ est une variable}  \\
\overline{s}(c)   =~& c    & \textrm{si $c$ est une constante} \\
\overline{s}([τ]) =~& [\overline{s}(τ)] & \\
\overline{s}(τ*)  =~& \overline{s}(τ)*   & \\
\overline{s}(τ_1, …, τ_n) \rightarrow τ_r =~& (\overline{s}(τ_1), …, \overline{s}(τ_n)) \rightarrow \overline{s}(τ_r) \\
\overline{s}(\{ f_1:τ_1, … ,f_n:τ_n \}) =~& \{ f_1:\overline{s}(τ_1), …, f_n:\overline{s}(τ_n) \}
\end{align*}

Par souci de simplicité, on notera $s$ pour $\overline{s}$.
\end{definition}

\begin{definition}[Instanciation]
Un schéma de types $σ = ∀ \vec{a} . τ$ peut être instancié en un type $μ$ s'il
existe une substition $s$ telle que :

\begin{itemize}
\item $\textrm{Dom}(s) ⊆ \vec{a}$
\item $s(τ) = μ$
\end{itemize}

On note alors $μ ≤ σ$.
\end{definition}

\begin{definition}[Variables libres]
Les variables libres d'un type sont l'ensemble des variables de types qui
apparaissent dans celui-ci :

\begin{align*}
\tFV{a}   =~& \{a\} & \textrm{si $a$ est une variable}  \\
\tFV{c}   =~& ∅     & \textrm{si $c$ est une constante} \\
\tFV{[τ]} =~& \tFV{τ} & \\
\tFV{τ*}  =~& \tFV{τ} & \\
\tFV{(τ_1, …, τ_n) \rightarrow τ_r} =~& \bigcup_{i=1}^n \tFV{τ_i} ∪ \tFV{τ_r} \\
\tFV{\{ f_1:τ_1, … ,f_n:τ_n \}} =~& \bigcup_{i=1}^n \tFV{τ_i}
\end{align*}

On étend cette définition aux schémas de types :

\[
\tFV{∀ \vec{a} . τ} = \tFV{τ} - \vec{a}
\]

ainsi qu'aux contextes de typage :

\begin{align*}
\tFV{ε} =& ∅ \\
\tFV{Γ, x:τ} =& \tFV{Γ} \cup \tFV{τ} \\
\tFV{Γ, f:σ} =& \tFV{Γ} \cup \tFV{σ}
\end{align*}
\end{definition}

\begin{definition}[Généralisation]
La généralisation consiste à construire un schéma de type à partir d'un type, en
quantifiant sur les variables libres :

\[
\textrm{Gen}(τ, Γ) = ∀ \vec{a} . τ
\quad \textrm{où}
\quad \vec{a} = \tFV{τ} - \tFV{Γ}
\]
\end{definition}

En associant un schéma de type $σ$ à une fonction $f$, on indique que la
fonction pourra être utilisée avec tout type $τ$ qui est une instanciation de
$σ$.

\subsection{Environnements de typage}

Chaque jugement de typage est effectué dans un environnement de typage $Γ$
particulier, qui contient le contexte nécessaire : ici, le type des fonctions et
variables du programme.

\gramlr{Environnements}{
\begin{align*}
Γ  \gramisa & (Γ_{fun}, Γ_{var}) & \textrm{Fonctions, variables}
\end{align*}

\begin{align*}
Γ_{fun} \gramisa & ε           & \textrm{Environnement vide}
\\      \gramor  & Γ_{fun},f:σ & \textrm{Ajout d'une fonction}
\end{align*}

\begin{align*}
Γ_{var} \gramisa & ε           & \textrm{Environnement vide}
\\      \gramor  & Γ_{var},v:τ & \textrm{Ajout d'une variable}
\end{align*}
}

Lorsque ce n'est pas ambigü, si $Γ=(Γ_{fun},Γ_{var})$ on notera les extensions
d'environnement $Γ,f:σ$ pour $((Γ_{fun},f:σ), Γ_{var})$ et $Γ,x:τ$ pour
$(Γ_{fun},(Γ_{var},x:τ))$. De même, on notera $(f, σ) ∈ Γ$ si $(f, σ) ∈
Γ_{fun}$, ou $(x, τ) ∈ Γ$ si $(x, τ) ∈ Γ_{var}$.

\subsection{Jugements de typage}

Un des principes du typage est d'associer à chaque expression un type, qui
décrit le genre des valeurs produites par l'évaluation de cette fonction.

\begin{definition}[Jugement de typage]
Un jugement de typage est de la forme $Γ ⊢ e : τ$ et se lit "sous $Γ$, $e$ est
typable en $τ$".
\end{definition}

Les instructions et blocs, au contraire, n'ont pas de type.

\begin{definition}[Bloc bien typé]
On note $Γ ⊢ i$ pour "sous $Γ$, $i$ est bien typé", c'est à dire que ces sous
expressions sont typables en accord avec le flot de données (par exemple, pour
que l'instruction $lv \leftarrow e$ soit bien typée sous $Γ$, il faut que les
types de $lv$ et de $e$ puissent avoir le même type sous $Γ$).
\end{definition}

Le cas des fonctions est particulier puisque celles-ci ont un schéma de types
qui leur est associé.

\begin{definition}[Fonction bien typée]
On note $Γ ⊢ f : σ$ le fait qu'une fonction $f$ est typable en un schéma $σ$ dans $Γ$.
\end{definition}

Enfin, la notion de programme bien typé est intrinsèque : elle se fait
indépendemment d'un environnement externe.

\begin{definition}[Programme bien typé]
Un programme $P$ est bien typé s'il existe un environnement $Γ$ permettant de bien
typer toutes les composantes (fonctions, globales et bloc d'initialisation) d'un
programme. On notera alors $⊢ P$.
\end{definition}

\subsection{Programme}

Au niveau global, un programme P est bien typé (noté $⊢ P$) s'il existe un
environnement $Γ = (\vec{σ}, \vec{τ})$ permettant de typer ses composantes (les fonctions, les
globales et le bloc d'initialisation).

\[
\irule{Prog}{
  (\vec{σ}, \vec{τ}) ⊢ \vec{f} : \vec{σ} \\
  (\vec{σ}, \vec{τ}) ⊢ \vec{x} : \vec{t} \\
  (\vec{σ}, \vec{τ}) ⊢ b
}{
  ⊢ (\vec{f}, \vec{x}, b)
}
\]

\subsection{Flot de contrôle}

Les règles suivantes permettent de définir les jugements $Γ ⊢ i$. De manière
générale, les instructions sont bien typées si leurs sous-instructions sont bien typées.

\begin{mathpar}
\irule{Pass}{
}{
  Γ ⊢ ε
}
\and
\irule{Seq}{
  Γ ⊢ s \\
  Γ ⊢ b
}{
  Γ ⊢ s;b
}
\and
\irule{While}{
  Γ ⊢ b
}{
  Γ ⊢ \npkWhile{b}
}
\and
\irule{Goto}{
}{
  Γ ⊢ \npkGoto{l}
}
\and
\irule{DoWith}{
  Γ ⊢ b
}{
  Γ ⊢ \npkDoWith{b}{l}
}
\end{mathpar}

Dans le cas de la conditionnelle, il est en plus nécessaire de vérifier que la
condition est un entier.

\begin{mathpar}
\irule{If}{
  Γ ⊢ e : \tInt \\
  Γ ⊢ i_t \\
  Γ ⊢ i_f
}{
  Γ ⊢ \npkIf{e}{i_t}{i_f}
}
\end{mathpar}

\subsection{Left values}

On associe à chaque left-value un type, qui est aussi le type des valeurs que
celle-ci peut contenir. Le cas des variables est direct : si un couple
(variable, type) est dans l'environnement de typage, la variable possède ce
type.

\begin{mathpar}
\irule{Lv-Var}{
  (v, τ) ∈ Γ
}{
  Γ ⊢ v : τ
}
\end{mathpar}

Si une expression a un type pointeur, en la déréférençant on obtient une valeur
du type pointé.

\begin{mathpar}
\irule{Lv-Deref}{
  Γ ⊢ e : τ*
}{
  Γ ⊢ *e : τ
}
\end{mathpar}

\todo{accès à un élément d'un type composite}

\begin{mathpar}
\irule{Lv-Field}{
  Γ ⊢ lv : τ_s \\
  (f, τ_f) ∈ τ_s
}{
  Γ ⊢ lv.f : τ_f
}
\and
\irule{Lv-Array}{
  Γ ⊢ lv : [τ] \\
  Γ ⊢ e : \tInt
}{
  Γ ⊢ lv[e] : τ
}
\end{mathpar}

\subsection{Expressions}

Les constantes ont leurs types habituels. Notons que le pointeur nul (\cNil) a
un type polymorphe.

\begin{mathpar}
\irule{Const-Int}{
}{
  Γ ⊢ n : \tInt
}
\and
\irule{Const-Float}{
}{
  Γ ⊢ f : \tFloat
}
\and
\irule{Const-Nil}{
}{
  Γ ⊢ \cNil : τ*
}
\end{mathpar}

Un certain nombre d'opérations est possible sur le type \tInt.

\begin{mathpar}
\irule{Op-Int}{
  \opsymb ∈ \{+,-,\times,/,\&,|,\opxor,\&\&,||,⋘,⋙\} \\
  Γ ⊢ e_1 : \tInt \\
  Γ ⊢ e_2 : \tInt
}{
  Γ ⊢ e_1~\opsymb~e_2 : \tInt
}
\end{mathpar}

De même sur \tFloat.

\begin{mathpar}
\irule{Op-Float}{
  \opsymb ∈ \{+.,-.,\times.,/.\} \\
  Γ ⊢ e_1 : \tFloat \\
  Γ ⊢ e_2 : \tFloat
}{
  Γ ⊢ e_1~\opsymb~e_2 : \tFloat
}
\end{mathpar}

Les opérateurs de comparaison peuvent s'appliquer à deux opérandes de types
"comparables". On introduit donc un jugement $\textsc{Comparable}(τ)$ qui est
vrai pour les types \tInt, \tFloat et pointeurs. Les comparaisons renvoient
alors un \tInt.

\begin{mathpar}
\irule{Cmp-Num}
  {τ ∈ \{ \tInt, \tFloat \}}
  {\textsc{Comparable}(τ)}
\and
\irule{Cmp-Ptr}
  { }
  {\textsc{Comparable}(τ*)}
\and
\irule{Op-Cmp}{
  \opsymb ∈ \{=,≠,≤,≥,<,>\} \\
  Γ ⊢ e_1 : τ \\
  Γ ⊢ e_2 : τ \\
  \textsc{Comparable}(τ)
}{
  Γ ⊢ e_1~\opsymb~e_2 : \tInt
}
\end{mathpar}

Les opérateurs unaires de négations "$-$" et "$-.$" s'appliquent respectivement
aux \tInt et aux \tFloat.

\begin{mathpar}
\irule{Unop-Minus-Int}{ Γ ⊢ e : \tInt }{ Γ ⊢ - e : \tInt }
\and
\irule{Unop-Minus-Float}{ Γ ⊢ e : \tFloat }{ Γ ⊢ -. e : \tFloat }
\end{mathpar}

Les opérateurs de négation unaires, en revanche, ne s'appliquent qu'aux
entiers.

\begin{mathpar}
\irule{Unop-Not}{
  \opsymb ∈ \{ \sim, ! \} \\
  Γ ⊢ e : \tInt
}{
  Γ ⊢ \opsymb~e : \tInt
}
\end{mathpar}

On peut prendre l'adresse d'une left-value, et obtenir un pointeur vers
celle-ci. Le type du résultat est un type pointeur vers le type de base.

\begin{mathpar}
\irule{AddrOf}{
  Γ ⊢ lv : τ
}{
  Γ ⊢ \&lv : τ*
}
\end{mathpar}

La règle concernant les pointeurs sur fonction est similaire, à ceci près qu'une
fonction a un schéma de type, qu'il faut instancier afin d'obtenir un type
pointeur.

\begin{mathpar}
\irule{AddrOfFun}{
  Γ ⊢ f : σ \\
  τ ≤ σ
}{
  Γ ⊢ \&f : τ
}
\end{mathpar}

\subsection{Fonctions}

Pour typer une fonction, on commence par ajouter ses paramètres et ses variables
locales dans l'environnement de typage, et on type la définition de la fonction.
Le type résultant est généralisé.

\[
\irule{Fun}{
  Γ' = Γ ⊕ \{ args(f) = \vec{τ}, loc(f) = \vec{τ_v} \} \\
  Γ' ⊢ body(f) \\
  Γ' ⊢ !ret : τ_r
}{
  Γ ⊢ f : Gen(\vec{τ} \rightarrow τ_r, Γ)
}
\]

\subsection{Instructions}

Une affectation est bien typée si elle faite entre une left-value et une
expression de même type.

\begin{mathpar}
\irule{Assign}{
  Γ ⊢ lv : τ \\
  Γ ⊢ e  : τ
}{
  Γ ⊢ lv \leftarrow e
}
\end{mathpar}

\wip{}

\begin{mathpar}
\irule{Fcall}{
  Γ ⊢ lv : τ_{ret} \\
  Γ ⊢ fe : σ \\
  Γ ⊢ \vec{e} : \vec{τ} \\
  (\vec{τ} \rightarrow τ_r) ≤ σ
}{
  Γ ⊢ lv \leftarrow fe(\vec{e})
}
\end{mathpar}

\section{Limitations}
\subsection{Programmes non typables}
\subsection{Incohérences}


\chapter{Analyse de provenance des pointeurs}

\label{cha:qualuser}
Dans le chapitre~\ref{cha:typbase}, nous avons vu comment ajouter un système de
types forts statiques à un langage impératif. Ici, nous étendons ce système afin
de lui ajouter des \emph{qualificateurs de type} qui décrivent l'origine des
données. Ils permettent de restreindre certaines opérations sensibles à des
expressions dont la valeur est sûre.

\section{Éditions et ajouts}

\section{Passage sur le cas d'étude}

\section{Propriété d'isolation mémoire}

Le déréférencement d'un pointeur dont la valeur est contrôlée par l'utilisateur
ne peut se faire qu'à travers une fonction qui vérifie la sûreté de celui-ci.


\chapter{Analyse de terminaison des chaînes C}

\label{cha:qualsz}
\section{But}
\section{Annotation de \texttt{string.h}}


\part{Expérimentation}

\chapter{Implantation}

\label{cha:implem}
Dans ce chapitre, nous décrivons la mise en œuvre des analyses statiques
précédentes. Nous commençons par un tour d'horizon des représentations
intermédiaires possibles, avant de décrire celle retenue : Newspeak. La chaîne
de compilation est explicitée, partant de C pour aller au langage impératif
décrit dans le chapitre~\ref{cha:lang}. Enfin, nous donnons les détails d'un
algorithme d'inférence de types à la Hindley-Milner, reposant sur l'unification
et le partage de références.

\section{Langages intermédiaires}

Le langage C \cite{KandR,AnsiC} a été conçu pour être une sorte d'assembleur
portable, permettant décrire du code indépendamment de l'architecture sur
laquelle il sera compilé. Historiquement, c'est il a permis de créer Unix, et
ainsi de nombreux logiciels bas niveau sont écrits en C. En particulier, il
existe des compilateurs de C vers les différents langages machine pour à peu
près toutes les architectures.

\begin{figure}
  \centering

  \begin{tikzpicture}
  \node[draw] (me) {Langage intermédiaire};
  \node[draw,  above left of=me, xshift=-4cm] (fe1) {Langage source 1};
  \node[draw,  below left of=me, xshift=-4cm] (fe2) {Langage source 2};
  \node[draw, above right of=me, xshift=4cm] (be1) {Langage destination 1};
  \node[draw, below right of=me, xshift=4cm] (be2) {Langage destination 2};

  \draw[->] (fe1) -- (me);
  \draw[->] (fe2) -- (me);

  \draw[->] (me) -- (be1);
  \draw[->] (me) -- (be2);

  \path (me) edge [loop below] node {Optimisations, analyses} (me);

  \coordinate (legbase) at ($ (me) + (0,-2cm) $);

  \draw[bigbrace] (me.south east |- legbase)
               to node [yshift=-5mm] {Middle-end}
                  (me.south west |- legbase);
  \draw[bigbrace] (me.south west |- legbase)
               to node [yshift=-5mm]{Front-end}
                  (fe2.south west |- legbase);
  \draw[bigbrace] (be2.south east |- legbase)
               to node [yshift=-5mm]{Back-end}
                  (me.south east |- legbase);
\end{tikzpicture}


  \caption{Décomposition d'un compilateur : front-ends, middle-end, back-ends}
  \label{fig:middle-end}
\end{figure}

Lors de l'écriture d'un compilateur, on a besoin d'un langage intermédiaire qui
fasse l'intermédiaire entre \emph{front-end} et \emph{back-end}
(figure~\ref{fig:middle-end}). Depuis ce langage on doit pouvoir exprimer des
transformations intermédiaires sur cette représentation (analyses sémantiques,
optimisations, etc), mais aussi compiler ce langage vers un langage machine.

L'idée de prendre C comme langage intermédiaire est très séduisante, mais
malheureusement sa sémantique est trop complexe et trop peu spécifiée. Il est
donc judicieux d'utiliser un langage plus simple à cet effet. Dans de nombreux
projets, des sous-ensembles de C ont été définis pour aller dans ce sens.

\subsubsection{Langages}

Les premiers candidats sont bien entendu les représentations intermédiaires
utilisées dans les compilateurs C. Elles ont l'avantage d'accepter en plus du C
standard, les diverses extensions (GNU, Microsoft, Plan9) utilisées par la
plupart des logiciels. En particulier, le noyau Linux repose fortement sur les
extensions GNU.

\paragraph{GCC} utilise une représentation interne nommée
GIMPLE\cite{gcc-gimple}. Il s'agit d'une structure d'arbre écrite en C, reposant
sur de nombreuses macros afin de cacher les détails d'implémentation pouvant
varier entre deux versions de GCC. Cette représentation étant réputée difficile
à manipuler, le projet MELT\cite{gcc-melt} permet de générer une passe de
compilation à partir d'un dialecte de Lisp.

\paragraph{LLVM}\cite{llvm-pres} est un compilateur développé par la communauté
puis sponsorisé Apple. À la différence de GCC, sa base de code est écrite en
C++. Il utilise une représentation intermédiaire qui peut être manipulée soit
sous forme d'une structure de données C++, soit d'un fichier de code-octet
compact, soit sous forme textuelle.

\paragraph{Objective Caml}\link{ocamlManual} utilise pour sa génération de code
une représentation interne nommée Cmm, disponible dans les sources du
compilateur sous le chemin \texttt{asmcomp/cmm.mli} (il s'agit donc d'une
structure de données OCaml). Ce langage a l'avantage d'être très restreint, mais
malheureusement il n'existe pas directement de traducteur permettant de compiler
C vers Cmm.

\paragraph{C--}\cite{spjcmm} \link{cmm}, dont le nom est inspiré du précédent,
est un projet qui visait à unifier les langages intermédiaires utilisés par les
compilateurs. L'idée est que si un front-end peut émettre du C-- (sous forme de
texte), il est possible d'obtenir du code machine efficace. Le compilateur
Haskell GHC utilise une représentation intermédiaire très similaire à C--.

Comme le problème de construire une représentation intermédiaire adaptée à une
analyse statique n'est pas nouveau, plusieurs projets ont déjà essayé d'y
apporter une solution. Puisque qu'ils sont développés en parallèle des
compilateurs, le support des extensions est en général moins important dans ces
langages.

\paragraph{CIL}\cite{NeculaCil} \link{kerneisCil} est une représentation en
OCaml d'un programme C, développée depuis 2002. Grâce à un mécanisme de
greffons, elle permet de prototyper rapidement des analyses statiques de
programmes.

\paragraph{Newspeak}\cite{newspeak} est un langage intermédiaire développé par
EADS Innovation Works, et qui est spécialisé dans l'analyse de valeurs par
interprétation abstraite. Il sera décrit plus en détails dans la
section~\ref{sec:npk}.

\paragraph{Compcert} est un projet qui vise à produire un compilateur certifié
pour C. C'est à dire que le fait que les transformations conservent la
sémantique est prouvé. Il utilise de nombreux langages intermédiaires, dont CIL.
Pour le front-end, le langage se nomme Clight\cite{cfront}. Les passes de
middle-end, quant à elles, utilisent Cminor\cite{cminorSL}.

%\begin{tabular}{| l || l | l | l | l | l | l |}
\hline
Nom           & Langage hôte    & Description       & But                 & Langage source                   &  Sémantique & Types \\ \hline \hline
C             & Texte           & \cite{AnsiC}      & Compilation         & C                                &  Non        & Oui   \\ \hline
GIMPLE        & C               & \cite{gcc-gimple} & Compilation         & C + extensions, C++, Ada, \ldots &  Non        & Non   \\ \hline
LLVM          & C++             & \cite{llvm-pres}  & Compilation         & C + extensions, C++, Ada, \ldots &             & Oui   \\ \hline
Cminusminus   & Texte           & \cite{spjcmm}     & Compilation         &                                  &             & Oui   \\ \hline
Cmm (Haskell) & Texte / Haskell &                   & Compilation         &                                  &             &       \\ \hline
Cmm (OCaml)   & OCaml           &                   & Compilation         &                                  &             &       \\ \hline
CIL           & OCaml           & \cite{NeculaCil}  & Analyse             &                                  &             &       \\ \hline
Clight        & Coq             & \cite{cfront}     & Compilation/Analyse &                                  &  Oui        & Oui   \\ \hline
Cminor        & Coq             & \cite{cminorSL}   & Compilation/Analyse &                                  &  Oui        & Oui   \\ \hline
Newspeak      & Ocaml           & \cite{newspeak}   & Analyse             &                                  &  Oui        & Oui   \\ \hline
\end{tabular}

% vim: textwidth=0


\section{Newspeak}
\label{sec:npk}

\wip{}

\section{Chaîne de compilation}

\begin{figure}
  \centering
  \begin{tikzpicture}\shorthandoff{!}
\tikzstyle{file}=[draw, shape=rectangle, node distance=2.2cm, minimum
height=1cm, shade, top color=white,
    bottom color=blue!50!black!20, draw=blue!40!black!60, very thick];

\node [file] (c1) {\textcolor{black}{.c}};
\node [file, below of=c1] (c2) {\textcolor{black}{.c}};
\node [node distance=2.2cm, below of=c2] (c3) {};

\node [file, minimum height=0, node distance=5mm, above of=c3,draw] (c3b) {.adb};
\node [file, minimum height=0, node distance=5mm, below of=c3,draw] (c3s) {.ads};

\path (c3b.north west) ++(-3mm,3mm) [draw,dotted] rectangle ($(c3s.south east)+(3mm,-3mm)$);

\node [below of=c3, node distance=2.2cm] (c4){};

\node [file, right of=c1] (cc1) {\textcolor{black}{.c}};
\node [file, right of=c2] (cc2) {\textcolor{black}{.c}};
\node [node distance=2.2cm, right of=c3] (cc3) {};

\node [below of=cc3, node distance=2.2cm](cc4){};

\draw[->] (c1) -- node[above] {{\tiny \ttfamily cpp}} (cc1);
\draw[->] (c2) -- (cc2);

\node [file, right of=cc1] (no1) {\textcolor{black}{.no}};
\node [file, right of=cc2] (no2) {\textcolor{black}{.no}};
\node [file, right of=cc3] (no3) {\textcolor{black}{.no}};

\node [below of=no3, node distance=2.2cm](no4){};

\draw[->] (cc1) -- node[above] {{\tiny \ttfamily c2newspeak -c}} (no1);
\draw[->] (cc2) --  (no2);

\draw[->] ($ (c3b.north east)!0.5!(c3s.south east) + (3mm,0) $) -- node[above] {\tiny \ttfamily ada2newspeak -c} (no3);


\node [file, right of=no2] (npk) {\textcolor{black}{.npk}};

\node [right of=no3, node distance=2.2cm](npk2){};
\node [right of=no4, node distance=2.2cm](npk3){};

\node[draw, ellipse, right of=no2, minimum height=3cm]{};

\draw[->] (no1) -- node {} (npk);
\draw[->] (no2) -- node[above] {{\tiny \ttfamily c2newspeak}} (npk);
\draw[->] (no3) -- node {} (npk);


\node [file, right of=npk, node distance=3cm, yshift=1cm](warn)
    {\textcolor{black}{\parbox{2cm}{\centering Programme \linebreak typé}}};
\node [right of=npk, node distance=3cm, yshift=-1cm](warnn)
    {\textcolor{black}{Erreurs}};

\node [right of=npk3, node distance=3cm](analyser2){};

\node [right of=npk, node distance=3cm] {\small ou};

\draw[->] (npk) -- node[above, yshift=2mm] {{\tiny \ttfamily ptrtype}} (warn);

\draw[->] (npk) -- (warnn);

\draw[->] (c4)   -- node[above, text depth=3pt] {\footnotesize Pr\'etraitement} (cc4);
\draw[->] (cc4)  -- node[above, text depth=3pt] {\footnotesize Compilation} (no4);
\draw[->] (no4)  -- node[above, text depth=3pt] {\footnotesize \'Edition de liens} (npk3);
\draw[->] (npk3) -- node[above, text depth=3pt] {\footnotesize Analyse} (analyser2);
\end{tikzpicture}

  \caption{Compilation depuis Newspeak}
  \label{fig:compil-npk}
\end{figure}

\todo{Mettre à jour la figure}

La compilation vers C est faite en trois étapes (figure~\ref{fig:compil-npk}) :
prétraitement du code source, compilation de C prétraité vers \newspeak{}, puis
compilation de \newspeak{} vers ce langage.

\subsection{Prétraitement}

\ctonewspeak{} travaillant uniquement sur du code prétraité (dans directives de
préprocesseur), la première étape consiste donc à faire passer le code par \cpp:
les macros sont développées, les constantes remplacées par leurs valeurs, les
commentaires supprimés, les fichiers d'en-tête inclus, etc.

\subsection{Compilation (levée des ambigüités)}

Cette passe est réalisée par l'utilitaire \ctonewspeak{}. L'essentiel de la
compilation consiste à mettre à plat les définition de types, et à simplifier le
flot de contrôle. C en effet propose de nombreuses constructions ambigües ou
redondantes.

Au contraire, \newspeak{} propose un nombre réduit de constructions. Rappelons
que le but de ce langage est de faciliter l'analyse statique : des constructions
orthogonales permettent donc d'éviter la duplication de règles sémantique, ou de
code lors de l'implémentation d'un analyseur.

Par exemple, plutôt que de fournir une boucle \emph{while}, une boucle
\emph{do/while} et une boucle for, \newspeak{} fournit une unique boucle
\npkWhile{}. La sortie de boucle est compilée vers un \npkGoto{}, qui est
toujours un saut vers l'avant (similaire à un "break" généralisé).

La sémantique de \newspeak{} et la traduction de C vers \newspeak{} sont
décrites dans \cite{newspeak}. En ce qui concerne l'élimination des sauts vers
l'arrière, on peut se référer à \cite{goto}.

\subsection{Annotations}

\newspeak{} a de nombreux avantages, mais pour une analyse par typage il est
trop bas niveau. Par exemple, dans le code suivant

\begin{Verbatim}[commandchars=\\\{\}]
\PY{k}{struct} \PY{n}{s} \PY{p}{\PYZob{}}
    \PY{k+kt}{int} \PY{n}{a}\PY{p}{;}
    \PY{k+kt}{int} \PY{n}{b}\PY{p}{;}
\PY{p}{\PYZcb{}}\PY{p}{;}

\PY{k+kt}{int} \PY{n+nf}{main}\PY{p}{(}\PY{k+kt}{void}\PY{p}{)}
\PY{p}{\PYZob{}}
    \PY{k}{struct} \PY{n}{s} \PY{n}{x}\PY{p}{;}
    \PY{k+kt}{int} \PY{n}{y}\PY{p}{[}\PY{l+m+mi}{10}\PY{p}{]}\PY{p}{;}
    \PY{n}{x}\PY{p}{.}\PY{n}{b} \PY{o}{=} \PY{l+m+mi}{1}\PY{p}{;}
    \PY{n}{y}\PY{p}{[}\PY{l+m+mi}{1}\PY{p}{]} \PY{o}{=} \PY{l+m+mi}{1}\PY{p}{;}
    \PY{k}{return} \PY{l+m+mi}{0}\PY{p}{;}
\PY{p}{\PYZcb{}}
\end{Verbatim}


\wip{}

\subsection{Implantation de l'algorithme de typage}

Commençons par étudier le cas du lambda-calcul simplement typé
(figure~\ref{fig:stlc}).

\begin{figure}

\gramlr{Termes}{
\begin{align*}
  t \gramisa & x              & \textrm{Variable} \\
    \gramor  & λ x . t        & \textrm{Abstraction} \\
    \gramor  & t              & \textrm{Application} \\
    \gramor  & n              & \textrm{Entier} \\
    \gramor  & f              & \textrm{Flottant} \\
    \gramor  & (t, t)         & \textrm{Couple} \\
    \gramor  & \textrm{fst}~t & \textrm{Projection gauche} \\
    \gramor  & \textrm{snd}~t & \textrm{Projection droite}
\end{align*}
}

\gramlr{Types}{
\begin{align*}
  τ \gramisa &  \tInt          & \textrm{Entier}\\
    \gramor  &  \tFloat        & \textrm{Flottant}\\
    \gramor  & τ \rightarrow τ & \textrm{Fonction} \\
    \gramor  & τ \times τ      & \textrm{Produit}
\end{align*}
}

\gramlr{Contextes}{
\begin{align*}
  Γ \gramisa & ε     & \textrm{Contexte vide}\\
    \gramor  & Γ,x:τ & \textrm{Extension}
\end{align*}
}

\gramlr{Règles}{
\begin{mathpar}
\irule{Int}{ }{Γ⊢n:\tInt}
\and
\irule{Float}{ }{Γ⊢f:\tFloat}
\and
\irule{Var}{x:τ∈Γ}{Γ⊢x:τ}
\and
\irule{App}{Γ⊢f:τ_1 \rightarrow τ_2 \\ Γ ⊢ x : τ_1}{Γ ⊢ f x : τ_2}
\and
\irule{Abs}{Γ,x:τ_1⊢ y : τ_2}{Γ⊢ λx.y:τ_1 \rightarrow τ_2}
\and
\irule{Proj-g}{Γ⊢x:τ_1\times τ_2}{Γ⊢\textrm{fst}~x:τ_1}
\and
\irule{Proj-d}{Γ⊢x:τ_1\times τ_2}{Γ⊢\textrm{snd}~x:τ_2}
\and
\irule{Tup}{Γ⊢x:τ_1 \\ Γ⊢y:τ_2}{Γ⊢(x,y):τ_1 \times τ_2}
\end{mathpar}
}


\caption{Lambda calcul simplement typé avec entiers, flottants et couples}
\label{fig:stlc}

\end{figure}

Prenons l'exemple de la fonction suivante\footnote{ On suppose que \texttt{plus}
est une fonction de l'environnement global qui a pour type $\tInt \rightarrow
\tInt \rightarrow \tInt$.} :

\[
f = λx.λy. \textrm{plus} (\textrm{plus} (\textrm{fst} x) (\textrm{snd} x)) y
\]

On voit que puisque \texttt{fst} et \texttt{snd} sont appliqués à \texttt{x}, ce
doit être un tuple. En outre on additionne ces deux composantes ensemble, donc
elles doivent être de type \tInt (et le résultat aussi). Par le même argument,
\texttt{y} doit aussi être de type \tInt. En conclusion, \texttt{x} est de type
$\tInt \times \tInt$ et \texttt{y} de type $\tInt$, donc f est de type $\tInt
\times \tInt \rightarrow \tInt \rightarrow \tInt$.

Mais comment faire pour implanter cette analyse ? En fait le système de types de
la figure~\ref{fig:stlc} a une propriété particulièrement intéressante : chaque
forme syntaxique (variable, abstraction, etc) est en conclusion exactement d'une
règle de typage. Cela permet de toujours savoir quelle règle il faut appliquer.

Partant du terme de conclusion ($f$), on peut donc en déduire un squelette
d'arbre d'inférence (figure~\ref{fig:inftree-rules})\footnote{Par souci de
clarté, les prémisses des applications de \textsc{(Var)} ne sont pas notées.}

\begin{figure} % {{{ Fig règles
\def\disptypeL#1{}
\def\disptypeR#1{}

\def\gammaXY{\ensuremath{Γ_2}}
\def\gammaX{\ensuremath{Γ_1}}
\def\gammaOneDef{\ensuremath{Γ^0, x\disptypeL{\tInt \times \tInt}, y\disptypeL{\tInt}}}
\def\gammaTwoDef{\ensuremath{Γ^0, x\disptypeL{\tInt \times \tInt}}}
\def\termFst#1{\ensuremath{\textrm{fst}~x}}
\def\termSnd#1{\ensuremath{\textrm{snd}~x}}
\def\fstx{\termFst{x}}
\def\sndx{\termSnd{x}}
\begin{mathpar}
\irule
  {Abs}
  {
    \irule
      {Abs}
      {
        \inferrule*[right=(App)]
          {
            \inferrule*[right=(Var), rightskip=4em]
              { }
              {\gammaXY ⊢ y
              \disptypeR{\tInt}
              }
            \\
            \inferrule*[right=(App),vdots=4em,leftskip=5em,rightskip=5em]
              {
                \inferrule*[right=(Var), leftskip=3em, rightskip=7em]
                  { }
                  {\gammaXY ⊢ \textrm{plus}
                  \disptypeR{\tInt \rightarrow \tInt \rightarrow \tInt}
                  }
                \\
                \inferrule*[right=(App),vdots=3em,leftskip=7em,rightskip=3em]
                  {
                    \inferrule*[right=(App), vdots=4em, rightskip=5em]
                      {
                        \inferrule*[right=(Var), rightskip=5em]
                          { }
                          {\gammaXY ⊢ \textrm{plus}
                          \disptypeR{\tInt \rightarrow \tInt \rightarrow \tInt}
                          }
                        \\
                        \inferrule*[right=(Proj-G), vdots=3em, leftskip=5em]
                          {
                            \irule
                              {Var}
                              { }
                              {\gammaXY ⊢ x
                              \disptypeR{\tInt \times \tInt}
                              }
                          }
                          {\gammaXY ⊢ \fstx
                          \disptypeR{\tInt}
                          }
                      }
                      {\gammaXY ⊢ \textrm{plus} (\fstx)
                      \disptypeR{\tInt \rightarrow \tInt}
                      }
                    \\
                    \inferrule*[right=(Proj-D), leftskip=5em, rightskip=4em]
                      {
                        \irule
                          {Var}
                          { }
                          {\gammaXY ⊢ x
                          \disptypeR{\tInt \times \tInt}
                          }
                      }
                      {\gammaXY ⊢ \sndx
                      \disptypeR{\tInt}
                      }
                  }
                  {\gammaXY ⊢ \textrm{plus} (\fstx) (\sndx)
                  \disptypeR{\tInt}
                  }
              }
              {\gammaXY ⊢ \textrm{plus} (\textrm{plus} (\fstx) (\sndx))
              \disptypeR{\tInt \rightarrow \tInt}
              }
          }
          {\gammaXY ⊢ \textrm{plus} (\textrm{plus} (\fstx) (\sndx)) y
          \disptypeR{\tInt}
          }
      }
      {\gammaX ⊢ λy. \textrm{plus} (\textrm{plus} (\fstx) (\sndx)) y
      \disptypeR{\tInt \rightarrow \tInt}
      }
  }
  {Γ_0 ⊢ λx.λy. \textrm{plus} (\textrm{plus} (\fstx) (\sndx)) y
  \disptypeR{\tInt \times \tInt \rightarrow \tInt \rightarrow \tInt}
  }
\end{mathpar}

\[\gammaX = \gammaTwoDef \qquad \gammaXY = \gammaOneDef\]


\caption{Arbre d'inférence : règles à utiliser}
\label{fig:inftree-rules}
\end{figure} % }}}

\begin{figure} % {{{ Fig full tree
\def\disptypeL#1{:#1}
\def\disptypeR#1{:#1}

\def\gammaXY{\ensuremath{Γ_2}}
\def\gammaX{\ensuremath{Γ_1}}
\def\gammaOneDef{\ensuremath{Γ^0, x\disptypeL{\tInt \times \tInt}, y\disptypeL{\tInt}}}
\def\gammaTwoDef{\ensuremath{Γ^0, x\disptypeL{\tInt \times \tInt}}}
\def\termFst#1{\ensuremath{\textrm{fst}~x}}
\def\termSnd#1{\ensuremath{\textrm{snd}~x}}
\def\fstx{\termFst{x}}
\def\sndx{\termSnd{x}}
\begin{mathpar}
\irule
  {Abs}
  {
    \irule
      {Abs}
      {
        \inferrule*[right=(App)]
          {
            \inferrule*[right=(Var), rightskip=4em]
              { }
              {\gammaXY ⊢ y
              \disptypeR{\tInt}
              }
            \\
            \inferrule*[right=(App),vdots=4em,leftskip=5em,rightskip=5em]
              {
                \inferrule*[right=(Var), leftskip=3em, rightskip=7em]
                  { }
                  {\gammaXY ⊢ \textrm{plus}
                  \disptypeR{\tInt \rightarrow \tInt \rightarrow \tInt}
                  }
                \\
                \inferrule*[right=(App),vdots=3em,leftskip=7em,rightskip=3em]
                  {
                    \inferrule*[right=(App), vdots=4em, rightskip=5em]
                      {
                        \inferrule*[right=(Var), rightskip=5em]
                          { }
                          {\gammaXY ⊢ \textrm{plus}
                          \disptypeR{\tInt \rightarrow \tInt \rightarrow \tInt}
                          }
                        \\
                        \inferrule*[right=(Proj-G), vdots=3em, leftskip=5em]
                          {
                            \irule
                              {Var}
                              { }
                              {\gammaXY ⊢ x
                              \disptypeR{\tInt \times \tInt}
                              }
                          }
                          {\gammaXY ⊢ \fstx
                          \disptypeR{\tInt}
                          }
                      }
                      {\gammaXY ⊢ \textrm{plus} (\fstx)
                      \disptypeR{\tInt \rightarrow \tInt}
                      }
                    \\
                    \inferrule*[right=(Proj-D), leftskip=5em, rightskip=4em]
                      {
                        \irule
                          {Var}
                          { }
                          {\gammaXY ⊢ x
                          \disptypeR{\tInt \times \tInt}
                          }
                      }
                      {\gammaXY ⊢ \sndx
                      \disptypeR{\tInt}
                      }
                  }
                  {\gammaXY ⊢ \textrm{plus} (\fstx) (\sndx)
                  \disptypeR{\tInt}
                  }
              }
              {\gammaXY ⊢ \textrm{plus} (\textrm{plus} (\fstx) (\sndx))
              \disptypeR{\tInt \rightarrow \tInt}
              }
          }
          {\gammaXY ⊢ \textrm{plus} (\textrm{plus} (\fstx) (\sndx)) y
          \disptypeR{\tInt}
          }
      }
      {\gammaX ⊢ λy. \textrm{plus} (\textrm{plus} (\fstx) (\sndx)) y
      \disptypeR{\tInt \rightarrow \tInt}
      }
  }
  {Γ_0 ⊢ λx.λy. \textrm{plus} (\textrm{plus} (\fstx) (\sndx)) y
  \disptypeR{\tInt \times \tInt \rightarrow \tInt \rightarrow \tInt}
  }
\end{mathpar}

\[\gammaX = \gammaTwoDef \qquad \gammaXY = \gammaOneDef\]


\caption{Arbre d'inférence complet}
\label{fig:inftree-full}
\end{figure} % }}}

Une fois à cette étape, on peut donner un nom à chaque type inconnu : $τ_1, τ_2,
\ldots$. L'utilisation qui en est faite permet de générer un ensemble de
contraintes d'unification. Par exemple, pour chaque application de la règle
\textsc{(App)} :

\[
\irule{App}{Γ⊢\ldots:τ_3 \\ Γ ⊢ \ldots : τ_1}{Γ ⊢ \ldots : τ_2}
\]

on doit déduire que $τ_3 = τ_1 \rightarrow τ_2$.

\clearpage

Ce signe $=$ est à prendre comme une contrainte d'égalité : partant d'un
ensemble de contraintes de la forme "type avec inconnue = type avec inconnue",
on veut obtenir une substitution "inconnue -> type concret".

Pour résoudre ces contraintes, on commence par les simplifier : si $τ_a
\rightarrow τ_b = τ_c \rightarrow τ_d$, alors $τ_a = τ_c$ et $τ_b = τ_d$. De
même si $τ_a \times τ_b = τ_c \times τ_d$. Au contraire, si $τ_a \rightarrow τ_b
= τ_c \times τ_d$, il est impossible d'unifier les types et il faut abandonner
l'inférence de types. D'autre cas sont impossibles, par exemple $\tInt = τ_1
\rightarrow τ_2$ ou $\tInt = \tFloat$.

Une fois ces simplifications réalisées, les contraintes restantes sont d'une des
formes suivantes :

\begin{itemize}
\item
  $τ_i = τ_i$. Il n'y a rien à faire, cette contrainte peut être supprimée.
\item
  $τ_i = τ_j$ avec $i \ne j$ : toutes les occurrences de $τ_j$ dans les autres
  contraintes peuvent être remplacées par $τ_i$.
\item
  $τ_i = x$ (ou $x = τ_i$) où $x$ est un type concret : idem.
\end{itemize}

\todo{C'est faux}

Une fois toutes les substitutions effectuées, on obtient un arbre de typage
correct (figure~\ref{fig:inftree-full}, donc un programme totalement inféré.

\clearpage

\begin{figure} % fig:unifpartage {{{

  \subfloat[][]{
  \label{fig:unifpartage:a}
  \begin{tikzpicture}
  \node               (var) {Var};
  \node[right of=var] (ref) {ref};
  \node[right of=ref, node distance=1.7cm] (u0) {Unknown 0};
  \node[below of=var] (ptr) {Ptr};
  \node[left of=var]  (x) {x :};
  \node[left of=ptr] (p) {p :};
  \draw[->] (x) -- (var);
  \draw[->] (p) -- (ptr);
  \draw[->] (ptr) -- (var);
  \draw[->] (var) -- (ref);
  \draw[->] (ref) -- (u0);
  \end{tikzpicture}
  }
  \subfloat[][]{
  \label{fig:unifpartage:b}
  \begin{tikzpicture}
  \node               (var) {Var};
  \node[right of=var] (ref) {ref};
  \node[right of=ref, node distance=1.7cm] (u0) {Instanciated};
  \node[right of=u0, node distance=1.7cm] (ii) {Int};
  \node[below of=var] (ptr) {Ptr};
  \node[left of=var]  (x) {x :};
  \node[left of=ptr] (p) {p :};
  \draw[->] (x) -- (var);
  \draw[->] (p) -- (ptr);
  \draw[->] (ptr) -- (var);
  \draw[->] (var) -- (ref);
  \draw[->] (ref) -- (u0);
  \draw[->] (u0) -- (ii);
  \end{tikzpicture}
  }
  \subfloat[][]{
  \label{fig:unifpartage:c}
  \begin{tikzpicture}
  \node               (var) {Int};
  \node[below of=var] (ptr) {Ptr};
  \node[left of=var]  (x) {x :};
  \node[left of=ptr] (p) {p :};
  \node[right of=ptr] (pi) {Int};
  \draw[->] (x) -- (var);
  \draw[->] (p) -- (ptr);
  \draw[->] (ptr) -- (pi);
  \end{tikzpicture}
  }

  \caption{Unification par partage}
  \label{fig:unifpartage}

\end{figure} % }}}

Plutôt que de modifier toutes les occurrences d'un type $τ_i$, on va affecter à
$τ_i$ la valeur du nouveau type.

L'implémentation de cet algorithme utilise le partage et les références
(figure~\ref{fig:unifpartage}).

D'abord \ref{fig:unifpartage:a}, ensuite \ref{fig:unifpartage:b}, et enfin
\ref{fig:unifpartage:c}.

\begin{figure} % fig:exunif:c {{{

  \insertcode{ex-unif-c.c}

  \caption{Compilation d'un programme C - avant}
  \label{fig:exunif:c}
\end{figure} % }}}

Prenons l'exemple de la figure~\ref{fig:exunif:c} et typons-le "à la main". On
commence par oublier toutes les étiquettes de type présentes dans le programme.
Celui-ci devient alors :

\begin{Verbatim}
var x, p;
p = &x;
x = 0;
\end{Verbatim}

La premiere ligne introduit deux variables. On peut noter leurs types respectifs
(inconnus pour le moment) $t_1$ et $t_2$. La première affectation \texttt{p =
\&x} implique que les deux côtés du signe "=" ont le même type. À gauche, le
type est $t_2$, et à droite $Ptr(t_1)$. On applique le même raisonnement à la
seconde affectation : à gauche, le type est $t_1$ et à droite Int. On en déduit
que le type de x est Int et celui de p est Ptr(Int).


\insertcode{lambda-types.ml}

Pour implanter cet algorithme, on représente les types de données du programmes
à typer par une valeur de type \texttt{ml\_type}. En plus des constantes de
types comme int ou float, et des constructeurs de type comme pair et fun, le
constructeur Var permet d'exprimer les variables de types (inconnues ou non).

Celles-ci sont numérotées par un int, on suppose avoir à disposition deux
fonctions manipulant un compteur global d'inconnues.

\begin{Verbatim}
module Counter : sig
  val reset_unknowns : unit -> unit
  val new_unknown : unit -> int
end
\end{Verbatim}

De plus, on a un module gérant les environnements de typage. Il pourra être
implanté avec des listes d'association ou des tables de hachage, par exemple. Sa
signature est :

\begin{Verbatim}
module Env : sig
  type t

  (* Construction *)
  val empty : t
  val extend : ml_ident -> ml_type -> t -> t

  (* Interrogation *)
  val get : ml_ident -> t -> ml_type option
end
\end{Verbatim}

Reprenons l'exemple précédent. Partant d'un environnement vide (Env.empty), on
commence par l'étendre de deux variables. Comme on n'a aucune information, il
fait allouer des nouveaux noms d'inconnues (qui correspondent à $t_1$ et $t_2$):

\begin{Verbatim}
let t1 = Var_type (Unknown (new_unknown ())) in
let t2 = Var_type (Unknown (new_unknown ())) in
let env =
  Env.extend "p" t2
    (Env.extend "x" t1
      Env.empty
    ) in
\end{Verbatim}

La première instruction indique que les deux côtés de l'affectation doivent
avoir le même type.

\begin{Verbatim}
let lhs1 = Lv_var "p"
and rhs1 = AddrOf (Exp_var "x") in
let t_lhs1 = typeof lhs1 env
and t_rhs1 = typeof rhs1 env in
unify t_lhs1 t_rhs1;
\end{Verbatim}

\begin{figure}
  \centering
  \begin{tikzpicture}
    [scale=2.3]

    \node at (0,  0) (n0a) {t1:};
    \node at (0, -1) (n0b) {t2:};

    \node at (1,  0) (n1a) {Var\_type};
    \node at (2,  0) (n2a) {ref};
    \node at (3,  0) (n3a) {Unknown 0};

    \node at (1, -1) (n1b) {Var\_type};
    \node at (2, -1) (n2b) {ref};
    \node at (3, -1) (n3b) {Instanciated};
    \node at (4, -1) (n4b) {Ptr};

    \draw [->] (n1a) -- (n2a);
    \draw [->] (n2a) -- (n3a);

    \draw [->] (n1b) -- (n2b);
    \draw [->] (n2b) -- (n3b);
    \draw [->] (n3b) -- (n4b);

    \draw [->] (n4b) -- ++ (0, 0.5) -| (n1a);

  \end{tikzpicture}
  \caption{Unification : partage}
  \label{fig:unifsharing}
\end{figure}

Ici il se passe plusieurs choses intéréssantes. D'une part nous faisont appel à
une fonction externe typeof qui retourne le type d'une expression sous un
environnement (dans une implantation complète il s'agirait d'un appel récursif).
Dans ce cas, \texttt{typeof lhs1 env} est identique à \texttt{Env.get lhs1 env}
et \texttt{typeof rhs1 env} à \texttt{Ptr\_type t1}. L'autre aspect intéressant
est la dernière ligne : la fonction \texttt{unify} va modifier en place les
représentations des types afin de les rendre égales. L'implantation de
\texttt{unify} sera décrite plus tard. Dans ce cas précis, elle va faire pointer
la référence dans t2 vers t1 (figure~\ref{fig:unifsharing}).

Enfin, la seconde affectation se déroule à peu près de la même manière.

\begin{Verbatim}
let lhs2 = Lv_deref (Lv_var "p")
and rhs2 = Exp_int 0 in
let t_lhs2 = typeof lhs2 env
and t_rhs2 = typeof rhs2 env in
unify t_lhs2 t_rhs2;
\end{Verbatim}

Ici \texttt{typeof lhs2 env} est identique à \texttt{Ptr\_type (Env.get "p"
env)} et \texttt{typeof lhs2 env} à \texttt{Const\_type Int\_type}. Et dans cas,
l'unification doit se faire entre t1 et \texttt{Const\_type Int\_type} : cela
mute la référence derrière t1 (figure~\ref{fig:typeunifref}).

\begin{figure}
  \begin{tikzpicture}
    [scale=2.3]

    \node at (0,  0) (n0a) {t1:};
    \node at (0, -1) (n0b) {t2:};

    \node at (1,  0) (n1a) {Var\_type};
    \node at (2,  0) (n2a) {ref};
    \node at (3,  0) (n3a) {Instanciated};
    \node at (4,  0) (n4a) {Const\_type};
    \node at (5,  0) (n5a) {Int\_type};

    \node at (1, -1) (n1b) {Var\_type};
    \node at (2, -1) (n2b) {ref};
    \node at (3, -1) (n3b) {Instanciated};
    \node at (4, -1) (n4b) {Ptr};

    \draw [->] (n1a) -- (n2a);
    \draw [->] (n2a) -- (n3a);
    \draw [->] (n3a) -- (n4a);
    \draw [->] (n4a) -- (n5a);

    \draw [->] (n1b) -- (n2b);
    \draw [->] (n2b) -- (n3b);
    \draw [->] (n3b) -- (n4b);

    \draw [->] (n4b) -- ++ (0, 0.5) -| (n1a);

  \end{tikzpicture}
  \caption{Unification par mutation de références}
  \label{fig:typeunifref}
\end{figure}

L'essence de l'algorithme d'inférence se situe donc dans 2 fonctions. D'une
part, \texttt{unify} qui réalise l'unification des types grâce à au partage des
références. D'autre part, la \texttt{typeof} qui encode les règles de typage
elles-mêmes et les applique à l'aide de \texttt{unify}.

\subsection{Algorithme d'unification}

Voici une implantation de la fonction \texttt{unify}.

Celle-ci prend en entrée deux types $t_1$ et $t_2$. À l'issue de l'exécution de
\texttt{unify}, ces deux types doivent pouvoir être considérés comme égaux. Si
ce n'est pas possible, une erreur sera levée.

La première étape est de réduire ces deux types, c'est à dire à transformer les
constructions \texttt{Var (ref (Instanciated t))} en \texttt{t}.

Ensuite, cela dépend des formes qu'ont les types réduits :

\begin{figure}
  \centering
  \begin{tikzpicture}
    [every node/.style={node distance=2cm}
    ,>=latex
    ]

    \node (n1) {Ptr};
    \node[right of=n1] (n2) {Var};
    \node[right of=n2] (n3) {ref};
    \node[right of=n3] (n4) {Instanciated};

    \draw [->] (n1) -- (n2);
    \draw [->] (n2) -- (n3);
    \draw [->] (n3) -- (n4);

    \draw [->] (n4) -- ++ (0, -1cm) -| (n1);

  \end{tikzpicture}
  \caption{Cycle dans le graphe de types}
  \label{fig:typecycle}
\end{figure}

\begin{itemize}

\item si les deux types sont inconnus (de la forme \texttt{Var (ref
(Instanciated t))}), on fait pointer l'une des deux références vers le premier
type. Notons que cela créé un type de la forme \texttt{Var (ref (Instanciated
(Var (ref (Unknown n)))))} qui sera réduit lors d'une prochaine étape
d'unification.

\item si un type est inconnu et pas l'autre, il faut de la même manière affecter la
référence. Mais en faisant ça inconditionnellement, cela peut poser problème :
par exemple en tentant d'unifier \texttt{a} avec \texttt{Ptr(a)} on pourrait
créer un cycle dans le graphe (figure~\ref{fig:typecycle}).
Pour éviter cette situation, il suffit de s'assurer que le type inconnu n'est
pas présent dans le type à affecter.

\item si les deux types sont des types de base (comme \tInt ou \tFloat), on ne
fait rien s'ils sont égaux ; l'algorithme échoue sinon.

\item si les deux types sont des constructeurs de type, il faut que les
constructeurs soient égaux. On unifie en outre leurs arguments deux à deux.

\item TODO sous typage pour les structures

\end{itemize}


TODO :

\begin{itemize}
\item implem du polymorphisme
\item implem du sous-typage
\item généralisation depuis le toy language
\end{itemize}

\clearpage

\begin{figure} % fig:exunif:tpk {{{

  \insertcode{ex-unif-tpk.ml}

  \caption{Compilation d'un programme C - après}
  \label{fig:exunif:tpk}
\end{figure} % }}}

Le programme C (figure~\ref{fig:exunif:c}) est compilé ainsi en Tyspeak
(figure~\ref{fig:exunif:tpk}).

% vim: spelllang=fr


\chapter{Étude de cas : un pilote de carte graphique}
\section{Description du problème}

\wip{}

% But
% What you're trying to tell the audience that they don't already know (e.g. Your story.)

Un système d'exploitation moderne comme GNU/Linux est séparé en deux niveaux de
privilèges : le noyau, qui gère directement le matériel, et les applications de
l'utilisateur, qui communiquent avec le noyau par l'interface restreinte des
\emph{appels système}.

Pour assurer l'isolation, ces deux parties n'ont pas accès aux mêmes zones
mémoire (cf. figure~\ref{fig:memmap}).

Si le code utilisateur tente d'accéder à la mémoire du noyau, une erreur sera
déclenchée. En revanche, si cette écriture est faite au sein de l'implantation
d'un appel système, il n'y aura pas d'erreur puisque le noyau a accès à toute la
mémoire : l'isolation aura donc été brisée.

Pour celui qui implante un appel système, il faut donc empêcher qu'un pointeur
passé en paramètre référence le noyau. Autrement dit, il est indispensable de
vérifier dynamiquement que la zone dans laquelle pointe le paramètre est
accessible par l'appelant\cite{hardy88confused}.

Si au contraire un tel pointeur est déréférencé sans vérification (avec
\texttt{*} ou une fonction comme \texttt{memcpy}), le code s'exécutera
correctement mais en rendant le système vulnérable, comme le montre la
figure~\ref{fig:radeon-bug}.

% Figure bug {{{

\begin{figure}
  cf annexe~\ref{cha:code-noyau}

  \caption[Bug freedesktop.org \#29340]{Bug freedesktop.org \#29340. Le
  paramètre \texttt{data} provient de l'espace utilisateur via un appel système.
  Un appelant malveillant peut se servir de cette fonction pour lire la mémoire
  du noyau à travers le message d'erreur.}

  \label{fig:radeon-bug}
\end{figure}
% }}}

Pour éviter cela, le noyau fournit un ensemble de fonctions qui permettent de
vérifier dynamiquement la valeur d'un pointeur avant de le déréférencer. Par
exemple, dans la figure précédente, la ligne 8 aurait dû être remplacée par :

\begin{Verbatim}
copy_from_user(&value, value_ptr, sizeof(value));
\end{Verbatim}

L'analyse présentée ici permet de vérifier automatiquement et statiquement que
les pointeurs qui proviennent de l'espace utilisateur ne sont déréférencés qu'à
travers une de ces fonctions sûres.

\section{Principes de l'analyse}

% Why the audience should believe that the results you've got aren't made up or
% flawed
Le problème est modélisé de la façon suivante : on associe à chaque variable
\texttt{x} un type de données \texttt{t}, ce que l'on note \texttt{x:t}. En
plus des types présents dans le langage C, on ajoute une distinction
supplémentaire pour les pointeurs. D'une part, les pointeurs ``noyau'' (de type
\texttt{t~*}) sont créés en prenant l'adresse d'un objet présent dans le code
source. D'autre part, les pointeurs ``utilisateurs'' (leur type est noté
\texttt{t user*}) proviennent des interfaces avec l'espace utilisateur.

Il est sûr de déréférencer un pointeur noyau, mais pas un pointeur
utilisateur. L'opérateur \texttt{*} prend donc un \texttt{t *} en entrée
et produit un \texttt{t}.

Pour faire la vérification de type sur le code du programme, on a besoin de
quelques règles. Tout d'abord, les types suivent le flot de données.
C'est-à-dire que si on trouve dans le code \texttt{a = b}, \texttt{a} et
\texttt{b} doivent avoir un type compatible. Ensuite, le qualificateur
\texttt{user} est récursif : si on a un pointeur utilisateur sur une structure,
tous les champs pointeurs de la structure sont également utilisateur. Enfin, le
déréférencement s'applique aux pointeurs noyau seulement : si le code contient
l'expression \texttt{*x}, alors il existe un type \texttt{t} tel que
\texttt{x:t*} et \texttt{*x:t}.

Appliquons ces règles à l'exemple de la figure \ref{fig:radeon-bug} : on suppose
que l'interface avec l'espace utilisateur a été correctement annotée. Cela
permet de déduire que \texttt{data:void user*}. En appliquant la première règle
à la ligne 6, on en déduit que \texttt{info:struct drm\_radeon\_info user*}
(comme en C, on peut toujours convertir de et vers un pointeur sur
\texttt{void}).

Pour déduire le type de \texttt{value\_ptr} dans la ligne 7, c'est la
deuxième règle qu'il faut appliquer : le champ \texttt{value} de
la structure est de type \texttt{uint32\_t~*} mais on y accède à travers
un pointeur utilisateur, donc \texttt{value\_ptr:uint32\_t user*}.

À la ligne 8, on peut appliquer la troisième règle : à cause du déréférencement,
on en déduit que \texttt{value\_ptr:t *}, ce qui est une contradiction puisque
d'après les lignes précédentes, \texttt{value\_ptr:uint32\_t user*}.

Si la ligne 3 était remplacée par l'appel à \texttt{copy\_from\_user}, il n'y
aurait pas d'erreur de typage car cette fonction peut accepter les arguments
\texttt{(uint32\_t~*, uint32\_t user*, size\_t)}.

\section{Implantation}

Une implantation est en cours. Le code source est d'abord prétraité par
\texttt{gcc -E} puis converti en Newspeak~\cite{newspeak}, un langage destiné à
l'analyse statique. Ce traducteur peut prendre en entrée tout le langage C, y
compris de nombreuses extensions GNU utilisées dans le noyau. En particulier,
l'exemple de la figure~\ref{fig:radeon-bug} peut être analysé.

À partir de cette représentation du programme et d'un ensemble d'annotations
globales, on propage les types dans les sous-expressions jusqu'aux feuilles.

Si aucune contradiction n'est trouvée, c'est que le code respecte la propriété
d'isolation. Sinon, cela peut signifier que le code n'est pas correct, ou bien
que le système de types n'est pas assez expressif pour le code en question.

Le prototype, disponible sur \link{penjili}, fera l'objet d'une démonstration.

\section{Conclusion}

% Recap of your story and its implications

Nous avons montré que le problème de la manipulation de pointeurs non sûrs peut
être traité avec une technique de typage. Elle est proche des analyses menées
dans CQual~\cite{pldi99} ou Sparse~\link{sparse}.

% Limitations
% Why someone might doubt your story and what you've done to get rid of as much
% doubt as possible.

Plusieurs limitations sont inhérentes à cette approche : notamment, la présence
d'unions ou de \emph{casts} entre entiers et pointeurs fait échouer l'analyse.

% Extensions

Le principe de cette technique (associer des types aux valeurs puis restreindre
les opérations sur certains types) peut être repris. Par exemple, si on définit
un type ``numéro de bloc'' comme étant un nouvel alias de \texttt{int}, on peut
considérer que multiplier deux telles valeurs est une erreur.


\chapter{Conclusion}

\section{Limitations}

\section{Perspectives}



\appendix

\chapter{Code du module noyau}

\label{cha:code-noyau}
On inclut ici le code analysé dans le chapitre~\ref{cha:etudedecas}. On inclut à
la suite le contexte nécessaire pour comprendre ce code.

{ \small
\insertcode{radeon-buggy.c}
}


\backmatter

%\chapterstar{Liste des figures}

\listoftodos

\clearpage

\listoffigures

\chapter{Références web}

\insertlinks

\bibliographystyle{alpha}
\bibliography{typ}
\end{document}
